\chapter{操作系统对象、路径与契约}

\section{从文件保存建立系统全景}

编辑器里点“保存”，表面上只是把一段字节写进名为 \code{data} 的文件，再告诉用户“保存成功”。这一句话看起来很小，但它实际包含三个承诺：用户给出的内容被系统接收了；旧文件不会因为半路失败被随便破坏；如果程序或机器随后出错，系统还能说明哪个版本可信。

如果机器上只有一个函数、一个设备和一个永远不会失败的写入动作，保存可以被想象成普通
函数调用。真实机器同时存在更多条件：设备可能暂时不能接收数据；用户传入的地址可能无效；
文件名和 \code{fd} 只是引用；写入可能先停在内存缓存里；断电后还要判断旧版本、新版本、
临时文件和目录项哪个可信。每一项条件都需要明确的状态和完成承诺，无法由一行函数返回值概括。

保存按钮可以作为一张 OS 地图。简单写法一旦说不清“谁在等、等什么、失败在哪里、成功凭什么”，就说明状态还藏在函数调用里，没有进入内核账本。能解释这些状态的对象，才是进程、内存、文件系统和驱动这些子系统真正要解决的问题。

\begin{keyidea}
操作系统的核心工作，是把“不可信的用户请求”和“会延迟、会失败的硬件事件”翻译成可管理的对象、可恢复的状态转移和可检查的完成证据。
\end{keyidea}

本书所说的“契约”指可检查的技术承诺。用户调用 \code{write} 时，OS 要说明怎样解释参数、
成功和失败分别表示什么；OS 提交设备请求后，硬件要用状态位、中断或完成队列报告请求是否
结束；文件系统宣布保存完成时，还要留下崩溃恢复能够识别的提交证据。任何一层缺少明确语义，
“保存成功”都无法成为可信结论。

\begin{longtable}{@{}L{0.16\linewidth}L{0.34\linewidth}L{0.42\linewidth}@{}}
\toprule
一句“保存成功” & 必须进一步证明什么 & 如果不证明会怎样 \\
\midrule
用户传入了内容 & 这段地址属于当前进程，且内核有权读取 & 可能读到非法地址，甚至泄漏内核数据 \\
系统接收了请求 & 当前执行流、文件引用和目标对象都被记录下来 & 设备慢时只能忙等，或完成后找不到发起者 \\
写入推进了 & 数据进入了哪个对象：内存缓存、设备队列还是持久介质 & \code{write} 返回后掉电，用户以为保存了，其实磁盘仍是旧内容 \\
失败能报告 & 错误发生在参数、权限、地址、设备、持久化还是目录发布 & 只能返回模糊失败，无法修复或恢复 \\
完成能证明 & 哪个提交点之后，崩溃恢复可以相信这个结果 & 重启后无法区分成功版本和半成品 \\
\bottomrule
\end{longtable}

这些名字不是孤立术语。\code{task} 记录当前执行流，VMA/PTE 解释用户地址，\code{fd} 和打开文件对象解释整数句柄，page cache 解释写入为何能先停在内存，DMA/completion 解释设备异步完成，\code{fsync} 解释怎样把内存状态推进到持久化边界。每个名字都对应保存路径上的一个缺口。

判断一个 OS 机制是否讲清楚，要看四类可观察证据：用户程序写了哪行代码；进入内核后新增了哪份状态账本；硬件用什么事件、寄存器或权限检查支撑这份账本；失败时外部能看到什么结果。比如 \code{write(fd, buf, 4)} 这一行，在用户层只是一个函数调用，在内核层要解释 fd、buffer 和长度，在硬件层可能触发 page fault、DMA 和中断，在证据层则只返回字节数或错误码。

先把这张地图画出来。本章后面的每一步修正，都会回到这张图上的某一层或某一个完成点。

\begin{pdfdiagram}
  \node[os compute, minimum width=64mm] (user) at (0,0) {用户程序\\\code{write(fd, buf, len)}};
  \node[os control, minimum width=64mm] (boundary) at (0,-14mm) {系统调用边界\\解释 fd、地址、长度，检查权限};
  \node[os state, minimum width=64mm] (objects) at (0,-28mm) {内核对象\\fd table、打开文件对象、VMA/PTE、page cache};
  \node[os memory, minimum width=64mm] (block) at (0,-42mm) {块层与驱动\\request、descriptor、doorbell、DMA};
  \node[os memory, minimum width=64mm] (device) at (0,-56mm) {设备与介质\\completion、flush、持久化};
  \draw[os strong bus] (user) -- (boundary);
  \draw[os strong bus] (boundary) -- (objects);
  \draw[os strong bus] (objects) -- (block);
  \draw[os strong bus] (block) -- (device);
  \node[os control, minimum width=52mm] (cp1) at (74mm,-14mm) {完成点 1\\\code{write} 返回：接口接受};
  \node[os control, minimum width=52mm] (cp2) at (74mm,-42mm) {完成点 2\\设备 completion：I/O 结束};
  \node[os control, minimum width=64mm] (cp3) at (32mm,-72mm) {完成点 3\\持久化发布：\code{fsync}、日志 commit、目录同步};
  \draw[os feedback] (boundary.east) -- (cp1.west);
  \draw[os feedback] (block.east) -- (cp2.west);
  \draw[os feedback] (device.south) -- (device.south |- cp3.north);
\end{pdfdiagram}
\diagramnote{保存请求自上而下穿过五层。三个“完成”在不同层生效：接口接受不等于设备完成，设备完成不等于崩溃后可恢复。掉电时能相信什么，取决于数据推进到了哪一层。}

把环境压到最小：系统中只有一个用户程序，只有一个用来写数据的设备，暂不考虑掉电，也没有别的任务抢 CPU。在这个受限场景里，保存过程可能写成这样：

\begin{lstlisting}
void save_request(const char* user_buf, size_t len) {
  while (!device_ready()) {
    // busy wait
  }

  for (size_t i = 0; i < len; ++i) {
    device_write_byte(user_buf[i]);
  }

  print("saved\n");
}
\end{lstlisting}

这段代码没有交代任何 OS 坐标：CPU 等待设备时谁拥有执行权，设备完成时谁负责唤醒，
\code{user\_buf} 是否可信，\code{fd} 指向什么，以及 \code{print("saved")} 依据哪项完成证据。
这些问题沿保存路径依次出现：CPU 先被无效占用，随后需要可靠睡眠和唤醒，接着要解释用户
buffer、fd 和文件对象，最后还要说明“保存成功”究竟越过了哪个边界。

先看等待设备的这一小段：

\begin{lstlisting}
while (!device_ready()) {
  // CPU keeps spinning here.
}
\end{lstlisting}

它反复读取同一个事实：设备现在能不能接收下一个字节。这个会影响下一步行为、必须被记住的事实，就是状态。

\begin{keyidea}
状态就是“会影响下一步行为的、必须被记住的事实”。在这里，\code{dev.ready == false} 表示设备还不能接收数据；\code{dev.ready == true} 表示保存路径可以继续往下走。
\end{keyidea}

忙等会让 CPU 持续询问同一个状态。假设设备 5 毫秒后才准备好，这段时间原本可以运行另一个
程序、处理网络包或响应输入；上面的 \code{while} 却把 CPU 固定在原地，使设备延迟直接
转化为处理器时间消耗。

\code{device\_ready()} 读取的是设备暴露给 CPU 的状态值。硬件设备常常会给 CPU 留出一些可读写的小格子，CPU 通过这些小格子知道设备是否空闲、是否出错、能不能接收下一个请求。这些小格子就是设备寄存器；在保存案例中，状态寄存器是 CPU 和设备之间交换状态的入口。

\begin{longtable}{@{}L{0.16\linewidth}L{0.22\linewidth}L{0.28\linewidth}L{0.26\linewidth}@{}}
\toprule
时刻 & 设备状态 & CPU 正在做什么 & 问题 \\
\midrule
0 ms & not ready & 读一次状态寄存器 & 发现不能写 \\
1 ms & not ready & 再读一次状态寄存器 & 没有任何新进展 \\
2 ms & not ready & 继续读状态寄存器 & 其他程序得不到 CPU \\
5 ms & ready & 终于跳出循环 & 前面几毫秒被浪费 \\
\bottomrule
\end{longtable}

因此，等待设备不能写成原地自旋。设备没准备好时，当前执行流应当先停下，把 CPU 让出来；设备准备好之后，内核再把它恢复到可运行状态。

要让执行流可靠地停下并恢复，需要引入 \code{task}。\code{task} 是内核能暂停、恢复、选择运行的一条执行流。它不是用户写的一个普通函数；它必须有一份账本，记录自己运行到哪里、现在能不能继续、如果不能继续是在等什么。

这里第一次出现“对象”。对象不是面向对象编程里的类名，也不是把名词包装得更正式。对象是内核为了保存某类状态而建立的账本。普通函数的局部变量只在当前调用栈里有效；一旦函数要睡眠、被中断、被调度走，或者等待另一个硬件事件回来，状态就不能只放在局部变量里。它必须放进一个能被调度器、中断处理路径、超时路径和关闭路径共同找到的地方。

\begin{longtable}{@{}L{0.18\linewidth}L{0.36\linewidth}L{0.38\linewidth}@{}}
\toprule
普通函数视角 & OS 对象视角 & 为什么必须升级 \\
\midrule
局部变量 \code{i} & 当前复制进度、请求偏移、已完成字节数 & 函数可能睡眠，醒来后要知道从哪里继续 \\
一次函数调用 & \code{task} 和保存的执行现场 & 当前执行流可能被切走，稍后从同一位置恢复 \\
一个布尔判断 & 谓词和受保护状态 & 条件可能被中断路径或另一个 CPU 修改 \\
等待一个结果 & wait queue 和 request/completion & 设备完成时要知道该唤醒谁、哪个请求完成 \\
\bottomrule
\end{longtable}

\begin{lstlisting}
if (!device_ready()) {
  current->state = BLOCKED;
  wait_queue_push(&dev.waiters, current);
  schedule();
}
\end{lstlisting}

这段伪代码引入了三个内核账本，分别解决“能不能运行”“等在哪里”“谁来选择下一个执行流”三个问题。

\begin{longtable}{@{}L{0.18\linewidth}L{0.35\linewidth}L{0.39\linewidth}@{}}
\toprule
词 & 在本例中的含义 & 反例 \\
\midrule
\code{BLOCKED} & 当前 task 不能继续运行，因为它在等设备 ready & 如果还把它放在可运行队列里，调度器可能又选中它，结果还是忙等 \\
\code{wait queue} & 等同一个条件的一组 task，例如都在等 \code{dev.ready} & 如果没有队列，设备完成时不知道该叫醒谁 \\
\code{schedule()} & 调度器入口：当前 task 让出 CPU，内核选择另一个能运行的 task & 如果没有调度器，睡眠只会变成整台机器停住 \\
\bottomrule
\end{longtable}

task 至少有三种基本状态：

\begin{longtable}{@{}L{0.18\linewidth}L{0.36\linewidth}L{0.38\linewidth}@{}}
\toprule
状态 & 含义 & 在保存例子里的位置 \\
\midrule
\code{running} & 正在某个 CPU 上执行 & 保存函数正在跑 \code{device\_ready()} 或复制数据 \\
\code{runnable} & 条件已经满足，只是暂时没轮到 CPU & 设备已经 ready，task 被唤醒，等调度器选中 \\
\code{blocked} & 现在不能推进，必须等某个条件变化 & 设备未 ready，task 睡在 \code{dev.waiters} 上 \\
\bottomrule
\end{longtable}

注意 \code{runnable} 不等于 \code{running}。一个 task 可以已经具备运行条件，但 CPU 正在跑别的 task；它只是排队等调度器。

\begin{pdfdiagram}
  \node[os compute, minimum width=30mm] (running) at (0,0) {running\\在 CPU 上执行};
  \node[os state, minimum width=30mm] (runnable) at (-34mm,-24mm) {runnable\\在 runqueue 里等 CPU};
  \node[os control, minimum width=30mm] (blocked) at (34mm,-24mm) {blocked\\在 wait queue 里等条件};
  \draw[os strong bus] (runnable) to[bend left=12] node[os label, left]{调度器选中} (running);
  \draw[os bus] (running) to[bend left=12] node[os label, right]{时间片到、被抢占} (runnable);
  \draw[os strong bus] (running) -- node[os label, right]{条件不满足，睡眠} (blocked);
  \draw[os strong bus] (blocked) -- node[os label, below]{条件变化，\code{wake\_up}} (runnable);
\end{pdfdiagram}
\diagramnote{调度器只在 running 和 runnable 之间做选择；blocked 的 task 不在候选集合里，必须先由唤醒路径把它挪回 runqueue。}

这三个教学状态在 Linux 的任务对象中有对应表示。任务状态字段记录能否运行，
队列成员关系再区分任务正在 CPU 上执行，还是已经进入就绪队列等待选择：

\begin{longtable}{@{}L{0.3\linewidth}L{0.28\linewidth}L{0.34\linewidth}@{}}
\toprule
真实状态名 & 对应教学状态 & 含义 \\
\midrule
\code{TASK\_RUNNING} & running 或 runnable & 在 CPU 上，或已在 runqueue 排队；两者共用一个状态名，靠 \code{on\_rq} 字段区分 \\
\code{TASK\_INTERRUPTIBLE} & blocked & 睡眠等待条件，可被信号打断 \\
\code{TASK\_UNINTERRUPTIBLE} & blocked & 睡眠等待条件（多为 I/O），信号打不断 \\
\code{\_\_TASK\_STOPPED} & blocked 的一种 & 被信号或调试器停住 \\
\code{EXIT\_ZOMBIE} & 已退出待回收 & 不再是执行候选，但保留退出状态等父进程读取 \\
\bottomrule
\end{longtable}

这个状态不需要相信教科书，可以直接在真实机器上看到。\code{/proc/<pid>/status} 的 \code{State} 字段就是 \code{\_\_state} 的用户态投影，下面两行是本书实验机上的真实输出：

\begin{lstlisting}
$ grep '^State' /proc/self/status
State:  R (running)
$ sleep 300 &
$ grep '^State' /proc/$!/status
State:  S (sleeping)
\end{lstlisting}

\code{R} 表示 \code{TASK\_RUNNING}——注意它同时覆盖“正在跑”和“在排队”两种情况，正好对应上面的状态机图；\code{S} 表示可中断睡眠，即 \code{TASK\_INTERRUPTIBLE}。后面章节的等待队列、唤醒和调度讨论，都会落到这几个真实状态名上。

仅把 task 放入等待队列还不够。前面的写法比忙等前进了一步，但它仍然不正确，问题出在“检查条件”和“登记等待者”之间：

\begin{lstlisting}
if (!device_ready()) {
  current->state = BLOCKED;
  wait_queue_push(&dev.waiters, current);
  schedule();
}
\end{lstlisting}

它把“检查条件”和“入等待队列”分成了几个互相可打断的动作。设备可能刚好在缝隙里完成：

\begin{longtable}{@{}L{0.12\linewidth}L{0.42\linewidth}L{0.38\linewidth}@{}}
\toprule
顺序 & 保存 task 做了什么 & 设备完成路径做了什么 \\
\midrule
1 & 检查 \code{dev.ready == false} & 还没发生 \\
2 & 准备把自己睡下去 & 设备变成 ready \\
3 & 还没入队 & 中断处理调用 \code{wake\_up(&dev.waiters)}，但队列为空 \\
4 & task 入队并 \code{schedule()} & 之后没有新的完成事件触发唤醒 \\
\bottomrule
\end{longtable}

这个错误叫 lost wakeup：唤醒确实发生过，但等待者还没登记，所以它错过了唤醒，之后可能永远睡下去。

这里的“唤醒”源于设备完成后交给 CPU 的事件，而非用户代码的一次普通函数调用。这类由设备
主动触发的事件称为中断。设备完成和用户函数继续执行分属不同路径，内核负责把前者转换为
“某个 task 已经具备继续运行的条件”。

lost wakeup 的根源是状态与等待者没有在同一规则下更新。等待者判断
\code{dev.ready == false} 时，完成路径可能马上把它改成 true；完成路径调用
\code{wake\_up} 时，等待者又可能尚未进入队列。两个动作虽然都发生了，却没有建立可靠关系。
因此，逐行检查代码还不够；不变量必须在所有可能的时序交错下成立。

\begin{pdfdiagram}
  \node[os title] at (-20mm,4mm) {等待路径（保存 task）};
  \node[os title] at (44mm,4mm) {完成路径（设备与中断）};
  \draw[os guide] (-20mm,1mm) -- (-20mm,-52mm);
  \draw[os guide] (44mm,1mm) -- (44mm,-52mm);
  \node[os small block, minimum width=46mm] (s1) at (-20mm,-10mm) {1 检查 \code{dev.ready == false}};
  \node[os small block, minimum width=46mm] (s2) at (44mm,-22mm) {2 设备完成，\code{ready = true}};
  \node[os small block, minimum width=46mm] (s3) at (44mm,-33mm) {3 \code{wake\_up}：队列还是空的};
  \node[os small block, minimum width=46mm] (s4) at (-20mm,-44mm) {4 入队并 \code{schedule}，无人再唤};
  \draw[os strong bus] (s1) -- (s4);
  \draw[os strong bus] (s2) -- (s3);
  \draw[os feedback] (s3.west) -- node[os label, above]{唤醒扑空} (s4.east |- s3.west);
  \draw[os feedback] (-52mm,-15mm) -- (-52mm,-39mm) node[os label, midway, left=1pt, align=right]{缝隙：检查之后\\入队之前};
\end{pdfdiagram}
\diagramnote{两条路径各自都“做了正确的事”，错误出在缝隙里：完成事件恰好落在等待者检查条件之后、登记入队之前。把检查和入队放进同一个受保护区间，缝隙才被封死。}

\begin{longtable}{@{}L{0.24\linewidth}L{0.33\linewidth}L{0.35\linewidth}@{}}
\toprule
错误理解 & 正确理解 & 后果 \\
\midrule
wakeup 是一条消息 & wakeup 只是提示条件可能变化 & 醒来后必须重新检查谓词 \\
入队和检查可以分开写 & 检查谓词和登记等待者必须受同一规则保护 & 否则完成事件会从缝隙里丢失 \\
睡眠只是暂停函数 & 睡眠会让出 CPU，并改变 task 的调度状态 & blocked task 不能继续留在 runqueue 上抢 CPU \\
锁只是防止变量同时写 & 锁保护的是“条件与队列的一致关系” & 没有锁会出现状态正确但等待关系错误 \\
\bottomrule
\end{longtable}

解决 lost wakeup，需要同时约束两个对象：一个是判断能否继续的条件，一个是保护这个条件和等待队列的互斥规则。

\begin{longtable}{@{}L{0.18\linewidth}L{0.36\linewidth}L{0.38\linewidth}@{}}
\toprule
词 & 在本例中是什么意思 & 为什么必须要它 \\
\midrule
谓词 & 一个能回答“现在能不能继续”的条件，这里是 \code{dev.ready} & wakeup 不是消息本身，task 醒来后还要重新检查条件 \\
锁 & 保护共享状态的一段互斥区间，这里保护 \code{dev.ready} 和 \code{dev.waiters} & 不能让完成路径插在“检查条件”和“入队”之间 \\
\bottomrule
\end{longtable}

可靠的等待路径要把“检查条件”和“入等待队列”放在同一个受保护区间里：

\begin{lstlisting}
lock(dev.lock);
while (!dev.ready) {
  prepare_to_wait(&dev.waiters, current);
  unlock(dev.lock);

  schedule();

  lock(dev.lock);
}
finish_wait(&dev.waiters, current);
unlock(dev.lock);
\end{lstlisting}

设备完成时，中断处理路径做相反的事：

\begin{lstlisting}
irq_handler() {
  lock(dev.lock);
  dev.ready = true;
  complete_current_request();
  wake_up(&dev.waiters);
  unlock(dev.lock);
}
\end{lstlisting}

这段代码有三个关键点。第一，用 \code{while}，不用 \code{if}，因为醒来只表示“条件可能变了”，不保证条件已经满足。第二，\code{prepare\_to\_wait} 发生在锁保护下，避免 lost wakeup。第三，\code{schedule()} 前要释放锁，否则设备完成路径拿不到锁，无法把 \code{dev.ready} 改成 true。

由此可以看到 OS 的核心形状：多条入口路径共同维护同一个谓词，正常执行只是这些路径形成的
一种状态序列。

\begin{longtable}{@{}L{0.2\linewidth}L{0.34\linewidth}L{0.38\linewidth}@{}}
\toprule
路径 & 它做什么 & 不能破坏的关系 \\
\midrule
等待路径 & 检查 \code{dev.ready}，不满足就把当前 task 放入 wait queue & 检查条件和入队不能被完成路径插到中间 \\
完成路径 & 设备中断到来，更新请求状态，再唤醒等待者 & 必须先写完成状态，再唤醒；否则等待者醒来仍看不到结果 \\
调度路径 & \code{schedule} 选择另一个 runnable task 运行 & blocked task 不能还留在 runqueue 里抢 CPU \\
返回路径 & 被唤醒后从 \code{schedule} 返回，重新检查条件 & wakeup 不是消息投递，只表示“条件可能变了” \\
\bottomrule
\end{longtable}

调度、锁、条件变量、futex 和 I/O completion 的共同骨架正是这组关系：先定义一个状态谓词，再定义谁能改它，最后用 wait queue 把“现在不能推进”的 task 挪开。

设备等待可靠之后，再处理 \code{user\_buf}。在普通 C 函数里，\code{const char* user\_buf} 看起来就是一个地址；但在 OS 里，它来自用户程序，内核不能直接相信。

具体地址如下：

\begin{lstlisting}
char* buf = (char*)0x1000;
write(fd, buf, 6000);
\end{lstlisting}

假设一页是 4096 字节，那么这次写入跨过两页：\code{0x1000..0x1fff} 和 \code{0x2000..0x276f}。第一页可能合法，第二页可能根本没有映射。如果内核把它当普通指针直接读，轻则内核自己 fault，重则把不该读的内存泄漏出去。

\begin{pdfdiagram}
  \node[os state, minimum width=36mm, minimum height=9mm] (vma) at (0,0) {};
  \node[os title] at (0,0) {VMA：声明过的可读区间};
  \node[os memory, minimum width=36mm] (p1) at (0,-13mm) {页 0x1000–0x1fff\\在 VMA 内，PTE present};
  \node[os control, minimum width=36mm] (p2) at (40mm,-13mm) {页 0x2000–0x2fff\\映射情况未知};
  \draw[os brace] (-18mm,-5mm) -- node[os label, above=2pt]{\code{user\_buf = 0x1000}，\code{len = 6000}} (55mm,-5mm);
  \node[os compute, minimum width=36mm] (d1) at (0,-31mm) {复制这 4096 字节};
  \node[os control, minimum width=36mm] (d2) at (40mm,-31mm) {page fault\\交给内核决策};
  \draw[os bus] (p1) -- (d1);
  \draw[os bus] (p2) -- (d2);
  \node[os label, anchor=west] at (62mm,-31mm) {在 VMA 内：补页后重试\\不在 VMA 内：\code{-EFAULT}};
\end{pdfdiagram}
\diagramnote{起点 \code{0x1000} 合法，只说明第一页可验证；buffer 是一个范围，检查必须按页推进到 \code{0x276f} 为止。第二页落入两种结局之一：合法但未装入（补页重试），或根本不属于任何 VMA（拒绝）。}

这里先引入系统调用边界。\code{write} 不是普通函数调用到底层设备；它会通过受控入口进入内核。用户态只交出三个值：fd、buffer 地址、长度。内核必须把这些不可信值变成自己能负责的状态。

最容易写错的检查是只看起点：

\begin{lstlisting}
if (is_user_pointer(user_buf)) {
  memcpy(kbuf, user_buf, len);
}
\end{lstlisting}

这段代码的问题和前面的忙等类似：它把一个连续过程压成了一个布尔判断。起点合法，只能说明第一个地址看起来属于用户空间；它不能证明后续 \code{len} 个字节都合法，也不能证明这些页现在都在内存里，更不能证明复制期间另一个线程不会改变映射。OS 需要的是按范围、按页、按权限推进，而不是对一个指针值投信任票。

\begin{longtable}{@{}L{0.22\linewidth}L{0.34\linewidth}L{0.36\linewidth}@{}}
\toprule
天真判断 & 漏掉的问题 & 正确追问 \\
\midrule
指针不为空 & 非空地址仍可能没有映射 & 它是否落在当前进程允许的地址区间内 \\
起点在用户空间 & 后续字节可能跨到非法页 & 整个范围是否逐页可读 \\
页表现在有翻译 & 复制中可能发生缺页、换入或权限变化 & fault 能否修复，修复后是否需要重试 \\
地址可读 & 目标对象未必允许写入 & fd、打开文件对象和权限还要单独检查 \\
\bottomrule
\end{longtable}

\begin{keyidea}
系统调用进入内核后，首先要解释用户传入的整数、指针和长度：fd 是否存在，地址是否可读，
长度计算是否越界，以及各类失败应返回什么错误。只有这些条件成立，实际操作才能开始。
\end{keyidea}

\code{user\_buf} 在 OS 里要拆成三层判断：

\begin{enumerate}
\tightlist
\item 这个地址是否落在当前进程声明过的某个地址区间里。
\item 这个区间是否允许当前操作，比如写文件时内核要从用户 buffer 读。
\item 当前页表里是否已经有可用翻译；没有时，缺页处理能不能补上。
\end{enumerate}

这三层判断落到四个概念上：

\begin{longtable}{@{}L{0.14\linewidth}L{0.42\linewidth}L{0.36\linewidth}@{}}
\toprule
词 & 含义 & 保存例子里的判断 \\
\midrule
虚拟地址 & 用户程序看到的地址编号，不直接等于物理内存位置 & \code{0x1000} 是用户给出的虚拟地址 \\
VMA & virtual memory area，一段地址区间的“使用意图” & 这段地址是不是用户栈、堆或 mmap 文件，是否允许读 \\
PTE & page table entry，一页地址的“当前翻译事实” & 这一页现在是否在内存里，是否允许用户态读 \\
page fault & CPU 发现地址翻译或权限不满足，把问题交给内核解释 & 第二页没有 PTE 时，内核决定是补页还是报错 \\
\bottomrule
\end{longtable}

因此，“读取用户 buffer”是一项受控复制过程，期间可能发生访问失败、缺页等待或部分完成。

这里的“可能睡眠”很关键。若第二页是合法文件映射，但内容还不在内存，内核可能需要从磁盘把页读进来；读磁盘又会提交块 I/O，当前 task 可能睡在缺页等待或 I/O completion 上。也就是说，一个看似单纯的 \code{copy\_from\_user}，可能临时走到调度器和块设备路径。系统调用不是一条不被打断的直线，它会在地址、权限、等待和设备完成之间来回穿行。

\begin{lstlisting}
size_t copied = 0;
while (copied < len) {
  int rc = copy_one_page_from_user(kbuf + copied,
                                   user_buf + copied);
  if (rc == PAGE_FAULT_FIXED) {
    copied += bytes_copied_this_round;
    continue;
  }
  if (rc == BAD_USER_ADDRESS) {
    break;
  }
}
\end{lstlisting}

两页复制过程如下：

\begin{longtable}{@{}L{0.18\linewidth}L{0.33\linewidth}L{0.41\linewidth}@{}}
\toprule
步骤 & 内核检查到什么 & 下一步 \\
\midrule
第一页 & 地址落在 VMA 内，VMA 允许读，PTE present & 复制第一页的数据 \\
第二页 & 地址落在 VMA 内，但 PTE not-present & 触发 page fault，能补页就重试 \\
第二页另一种情况 & 找不到覆盖它的 VMA & 停止复制，返回部分写入或 \code{-EFAULT} \\
内核地址 & PTE 标着“只允许内核访问”，用户无权访问 & 拒绝，不能泄漏内核内存 \\
\bottomrule
\end{longtable}

\begin{longtable}{@{}L{0.24\linewidth}L{0.34\linewidth}L{0.34\linewidth}@{}}
\toprule
用户传入的情况 & 内核看到什么 & 合理结果 \\
\midrule
buffer 全部合法且页已在内存 & VMA 允许读，PTE present & 复制成功，继续写入路径 \\
buffer 合法但页未装入 & VMA 允许读，PTE not-present & 触发 page fault，分配或读入页面后重试 \\
buffer 前半段合法，后半段非法 & 前几页复制成功，后一页找不到 VMA & 返回已写字节数或 \code{-EFAULT}，不能越界读 \\
buffer 指向内核地址 & 用户态无权访问内核专用页 & 拒绝，不能把内核内存泄漏给用户 \\
另一个线程同时修改 buffer & 指针仍合法，但内容在变化 & 内核复制到的是某一时刻的快照，不能假设之后不变 \\
\bottomrule
\end{longtable}

这个边界可以直接观察。下面的程序故意把 buffer 指向 \code{0x1}。这个地址通常不属于进程的有效 VMA；用户态把它作为普通整数传给内核并不违法，但内核不能相信它真的可读。

\begin{lstlisting}
int fd = open("data", O_WRONLY | O_CREAT | O_TRUNC, 0644);
const char* bad = (const char*)0x1;
ssize_t n = write(fd, bad, 16);
printf("n=%zd errno=%d\n", n, errno);
\end{lstlisting}

系统调用记录中的具体 fd 号只服务当前例子，更重要的信息是 \code{write} 怎样把坏地址转换成
用户可见错误：

\begin{lstlisting}
$ strace -e trace=write ./badbuf
write(3, 0x1, 16) = -1 EFAULT (Bad address)
\end{lstlisting}

程序自己看到的返回值同样能证明这条受控路径。本书实验机上的真实输出：

\begin{lstlisting}
$ ./badbuf
n=-1 errno=14 (Bad address)
\end{lstlisting}

这里没有出现“内核随便崩溃”，也没有出现“从地址 \code{0x1} 读出垃圾”。发生的是一条受控错误路径：CPU 在内核复制用户数据时发现地址访问不成立，进入 page fault；内核发现这次 fault 发生在 \code{copy\_from\_user} 这类允许恢复的位置，于是把它翻译成 \code{-EFAULT} 返回给系统调用。内核里通常会为这类可恢复访问准备异常表。异常表可以理解成一份小账本：如果某条特殊的内核访存指令 fault，不要按普通内核 bug 处理，而是跳到修复路径，设置错误结果并返回。

\begin{longtable}{@{}L{0.2\linewidth}L{0.34\linewidth}L{0.38\linewidth}@{}}
\toprule
观察到的现象 & 背后的边界 & 说明 \\
\midrule
\code{0x1} 能作为参数传入 & 系统调用入口只接收整数值 & 入口接收参数不代表参数语义有效 \\
\code{write} 返回 \code{EFAULT} & 地址复制边界失败 & 错误发生在读取用户 buffer，而不是 fd 或磁盘 \\
进程继续运行 & fault 被内核修复路径接住 & 这类 fault 是可报告错误，不是必须 panic 的内核 bug \\
文件不应被当作完整写入 & 复制没有得到 16 个可信字节 & 没有可信输入，就不能推进保存语义 \\
\bottomrule
\end{longtable}

这里有两个边界需要同时成立。第一，系统调用边界会把不可信输入变成内核拥有的稳定状态；小数据通常复制，大 I/O 可能 pin page，也就是临时固定页面，保证设备 DMA 完成前页面不会被回收或移动。第二，page fault 不是单纯的“内存不够”，它是硬件把一个地址访问问题交给内核解释：能修就修，不能修就返回错误或发信号。

\code{user\_buf} 解释完，还剩 \code{fd}。用户程序通常这样写：

\begin{lstlisting}
int fd = open("data.tmp", O_WRONLY | O_CREAT | O_TRUNC, 0644);
write(fd, buf, len);
\end{lstlisting}

\code{fd} 看起来只是一个整数，例如 3。但它不是文件本体。它是当前进程 fd table 里的一个索引。fd table 可以先理解成“这个进程手里的资源号码表”。fd table 指向的 file object 可以理解为“打开文件对象”：它表示一次 \code{open} 或 \code{dup} 关系背后的内核对象，可对应 POSIX 语义里的 open file description。

\begin{longtable}{@{}L{0.2\linewidth}L{0.34\linewidth}L{0.38\linewidth}@{}}
\toprule
层次 & 在本例中的含义 & 它回答的问题 \\
\midrule
fd 整数 & 用户态拿到的号码，例如 3 & 用户后续用哪个号码引用这次打开得到的对象 \\
fd table & 当前进程的号码表 & 3 号槽是否存在，指向哪个内核对象 \\
打开文件对象（file object） & 一次打开文件形成的内核对象 & 当前 offset、读写模式、引用计数、具体操作函数 \\
inode & 文件系统里的文件身份和元数据 & 这个文件的权限、大小、块映射、page cache 归属 \\
\bottomrule
\end{longtable}

\begin{pdfdiagram}
  \node[os compute, minimum width=28mm] (fdt) at (0,0) {fd table\\（每进程一张）};
  \node[os state, minimum width=44mm] (fo) at (52mm,0) {打开文件对象\\offset、模式、引用计数、操作表};
  \node[os memory, minimum width=42mm] (ino) at (108mm,0) {inode\\权限、大小、块映射、page cache};
  \draw[os strong bus] (fdt) -- node[os label, above]{fd 3 只是索引} (fo);
  \draw[os strong bus] (fo) -- node[os label, above]{指向文件身份} (ino);
  \node[os small block, minimum width=28mm] (fd3) at (0,-15mm) {fd 3};
  \node[os small block, minimum width=28mm] (fd4) at (0,-27mm) {fd 4（\code{dup}）};
  \draw[os bus] (fd3.east) -- (fo.south west);
  \draw[os bus] (fd4.east) -- (fo.south west);
  \node[os label, anchor=west] at (26mm,-33mm) {\code{dup}：两个 fd 槽共享同一个打开文件对象和 offset};
\end{pdfdiagram}
\diagramnote{整数 fd 只是进程 fd table 的索引；offset 和引用计数属于打开文件对象；inode 才是文件身份。\code{dup} 复制的是槽位，不是对象——这是后面所有共享 offset 现象的根源。}

为什么要分这么多层？一个 \code{dup} 例子能直接暴露 fd 和文件本体之间的差别：

\begin{lstlisting}
int a = open("data", O_WRONLY);
int b = dup(a);
write(a, "A", 1);
write(b, "B", 1);
\end{lstlisting}

\code{dup(a)} 会创建另一个 fd，但它们通常共享同一个打开文件对象，也就是共享这个对象里的 offset。因此第一次写入把 offset 从 0 推到 1，第二次写入接着从 1 写，文件内容会是 \code{AB}。如果把 fd 误认为“文件本体”，就解释不了为什么两个整数会共享 offset，也解释不了 \code{close(a)} 后 \code{b} 还能继续写。

这个共享 offset 可以通过一个小程序清楚地观察到：

\begin{lstlisting}
int a = open("data", O_WRONLY | O_CREAT | O_TRUNC, 0644);
int b = dup(a);

write(a, "A", 1);        // offset: 0 -> 1
write(b, "B", 1);        // offset: 1 -> 2
lseek(a, 0, SEEK_SET);   // shared offset: 2 -> 0
write(b, "C", 1);        // offset: 0 -> 1
\end{lstlisting}

如果 \code{a} 和 \code{b} 是两个互不相关的文件对象，\code{lseek(a, 0, SEEK_SET)} 不应该影响 \code{b}。但 \code{dup} 的结果会共享打开文件对象，所以 \code{b} 的下一次写入从 0 开始，覆盖第一个字节。最终文件内容不是 \code{ABC}，而是 \code{CB}。

\begin{lstlisting}
$ ./dup-offset
$ od -An -c data
   C   B
\end{lstlisting}

\begin{longtable}{@{}L{0.16\linewidth}L{0.28\linewidth}L{0.48\linewidth}@{}}
\toprule
动作 & 共享 offset & 文件内容 \\
\midrule
\code{write(a,"A",1)} & 0 变成 1 & \code{A} \\
\code{write(b,"B",1)} & 1 变成 2 & \code{AB} \\
\code{lseek(a,0)} & 2 变成 0 & \code{AB}，只是位置变了 \\
\code{write(b,"C",1)} & 0 变成 1 & \code{CB}，第一个字节被覆盖 \\
\bottomrule
\end{longtable}

这比口头说“fd 是引用”更有力。用户手里确实有两个整数，但内核账本里只有一个共享的打开文件对象；offset 属于这个对象，不属于整数变量本身。若代码需要两个互不影响的文件位置，就不能用 \code{dup} 得到第二个位置，而要重新 \code{open}，让内核创建另一个打开文件对象。

再看一个更容易在工程里踩到的例子：

\begin{lstlisting}
int fd = open("data", O_WRONLY);
int copy = dup(fd);
close(fd);
write(copy, "X", 1);
\end{lstlisting}

\code{close(fd)} 关闭的是当前进程 fd table 里的一个槽位；只要 \code{copy} 还引用同一个打开文件对象，这个打开对象就不能释放。于是“文件是否还能写”不取决于整数 \code{fd} 这个变量是否还存在，而取决于打开文件对象的引用计数和权限状态。在 \code{fork}、\code{exec}、socket 和设备文件中也会反复遇到同一条规则：用户态整数只是引用入口，真正的生命周期在内核对象上。

\begin{longtable}{@{}L{0.2\linewidth}L{0.35\linewidth}L{0.37\linewidth}@{}}
\toprule
现象 & 如果把 fd 当文件本体会怎样误解 & 用对象账本怎样解释 \\
\midrule
\code{dup} 后共享 offset & 以为两个整数各写各的 & 两个 fd 槽指向同一个打开文件对象 \\
\code{close(a)} 后 \code{b} 还能写 & 以为文件已经关闭 & 只减少一个引用，打开对象仍被 \code{b} 持有 \\
\code{fork} 后子进程继承 fd & 以为复制了整个文件 & 父子 fd table 槽可引用同一打开对象 \\
\code{exec} 后部分 fd 消失 & 以为程序替换会关闭全部资源 & \code{FD\_CLOEXEC} 决定哪些引用在提交点被裁剪 \\
\bottomrule
\end{longtable}

这层关系在真实系统里可以直接看到。下面的程序打开 \code{data.tmp}，用 \code{dup} 复制一个 fd，再打开所在目录，然后睡眠一分钟；趁它睡眠时去读它的 fd 目录（输出删去了完整路径）：

\begin{lstlisting}
$ ./fd_demo &
pid=26635 fd=3 copy=4 dirfd=5
$ ls -l /proc/26635/fd
3 -> .../data.tmp
4 -> .../data.tmp
5 -> .../
$ grep -E '^(State|Threads)' /proc/26635/status
State:  S (sleeping)
Threads:        1
\end{lstlisting}

fd 3 和 fd 4 指向同一个文件，这正是 \code{dup} 共享打开文件对象的直接证据；fd 5 是一个目录——目录也能被打开并占有 fd 槽位，后面的 \code{fsync(dirfd)} 依赖的正是这个事实。睡眠中的进程 \code{State} 是 \code{S}，对应前面讲的 \code{TASK\_INTERRUPTIBLE}。

这三层在 Linux 中分别由 fd 表、打开文件对象和 inode 对象表示。fd 表保存整数到对象引用的映射；
打开文件对象保存当前 offset、引用计数和操作方法；inode 表示持久文件身份，page cache 则与
该文件的映射对象关联。这里先确定一个结论：整数 fd、打开文件对象和 inode 是三层不同对象，
offset 属于中间的打开文件对象。

为了避免把 \code{dup} 的结论误记成“同一个路径就共享 offset”，再看重新 \code{open} 的反例：

\begin{lstlisting}
int a = open("data", O_WRONLY | O_CREAT | O_TRUNC, 0644);
int b = open("data", O_WRONLY);

write(a, "A", 1);        // a offset: 0 -> 1
write(b, "B", 1);        // b offset: 0 -> 1
\end{lstlisting}

这次最终内容通常是 \code{B}，不是 \code{AB}。原因不是路径名变了，而是两次 \code{open} 创建了两个打开文件对象；它们都指向同一个 inode，但各自有独立 offset。第一次写入用 \code{a} 的 offset 0，把文件变成 \code{A}；第二次写入用 \code{b} 的 offset 0，又从文件开头覆盖成 \code{B}。

\begin{lstlisting}
$ ./open-twice-offset
$ od -An -c data
   B
\end{lstlisting}

\begin{longtable}{@{}L{0.16\linewidth}L{0.39\linewidth}L{0.37\linewidth}@{}}
\toprule
写法 & 内核对象关系 & offset 结果 \\
\midrule
\code{dup(a)} & 两个 fd 槽指向同一个打开文件对象 & \code{a} 写完后，\code{b} 接着往后写 \\
重新 \code{open} & 两个打开文件对象指向同一个 inode & \code{a} 的 offset 不影响 \code{b} 的 offset \\
\code{fork} 继承 & 父子进程的 fd 槽可指向同一个打开文件对象 & 子进程推进 offset，父进程随后也从新位置写 \\
\bottomrule
\end{longtable}

\code{fork} 的情况也可以直接观察。下面的程序让子进程先写，父进程等子进程退出后再写：

\begin{lstlisting}
int fd = open("data", O_WRONLY | O_CREAT | O_TRUNC, 0644);
pid_t pid = fork();

if (pid == 0) {
  write(fd, "C", 1);
  _exit(0);
}

waitpid(pid, NULL, 0);
write(fd, "P", 1);
\end{lstlisting}

输出会显示父进程不是从 offset 0 覆盖子进程写入的字节，而是接在后面写：

\begin{lstlisting}
$ ./fork-offset
$ od -An -c data
   C   P
\end{lstlisting}

\code{waitpid} 只是在这个例子里固定执行顺序，真正决定 \code{CP} 的是父子进程共享同一个打开文件对象。子进程把共享 offset 从 0 推到 1；父进程被唤醒后继续使用同一个对象，所以从 1 开始写。由此可以把 fd 生命周期说清楚：进程可以复制，fd 槽可以复制，路径名可以相同，但打开文件对象才是 offset、状态标志和引用计数的承载者。

有了 fd 和 buffer 两条账，\code{write(fd, buf, len)} 的路径可以写成：

\begin{lstlisting}
write(fd, buf, len)
  -> find fd in current process fd table
  -> check the open file object allows write
  -> copy data from user buffer
  -> update file/page-cache state
  -> return byte count or -errno
\end{lstlisting}

这一阶段会用到 \code{page cache}。它是内核用内存缓存文件内容的一组页面。写普通文件时，内核常常先把用户数据复制到 page cache，把相关页面标记为 dirty，然后 \code{write} 就可以返回。dirty 的意思是“内存里的文件页已经比磁盘上的版本更新，还需要写回”。

page cache 保存的是文件内容在内核内存中的副本。数据进入 page cache 后，后续读取可以直接
复用，多个小写也能够合并成更合适的块 I/O；与此同时，完成边界增加了。\code{write} 返回时，
数据可能只到达内核内存；后台写回或 \code{fsync} 随后把它组织成块请求；completion 表示
相应设备请求结束；设备缓存和文件系统日志还会继续影响崩溃后哪些状态可信。

\begin{longtable}{@{}L{0.22\linewidth}L{0.34\linewidth}L{0.36\linewidth}@{}}
\toprule
边界 & 此时数据在哪里 & 还不能承诺什么 \\
\midrule
\code{copy\_from\_user} 结束 & 内核 buffer 或 page cache 中 & 目标文件名未必已经安全发布 \\
\code{write} 返回 & 接口语义接受了若干字节 & 设备未必见过这批数据 \\
块请求提交 & request 和 DMA 映射交给驱动/设备 & completion 未到，buffer 不能释放 \\
设备 completion & 某次设备请求完成或失败 & 介质持久性仍可能依赖 flush/FUA/日志 \\
\code{fsync} 成功 & 文件相关数据和必要元数据推进到持久化边界 & 目录项更新仍可能需要目录 \code{fsync} \\
\bottomrule
\end{longtable}

最后回到最初的 \code{print("saved")}。这行太早，因为保存请求至少有三个不同的完成点：

\begin{longtable}{@{}L{0.22\linewidth}L{0.34\linewidth}L{0.36\linewidth}@{}}
\toprule
完成点 & 说明 & 反例 \\
\midrule
系统调用完成 & \code{write} 返回，内核已经按接口语义接受了一部分数据 & 机器立刻掉电，dirty page 还没写到设备 \\
设备请求完成 & 块层或驱动收到 completion，某次 I/O 已结束 & 设备缓存还没 flush，断电后介质上未必有数据 \\
持久化发布完成 & \code{fsync}、日志 commit、\code{rename} 和目录同步形成可恢复状态 & 只同步文件内容，不同步目录项，崩溃后新名字可能丢失 \\
\bottomrule
\end{longtable}

\begin{pdfdiagram}
  \draw[os strong bus] (0,0) -- (134mm,0);
  \node[os small block, minimum width=24mm] (t0) at (0,11mm) {\code{write} 返回};
  \node[os small block, minimum width=26mm] (t1) at (34mm,11mm) {dirty page\\等在 page cache};
  \node[os small block, minimum width=26mm] (t2) at (67mm,11mm) {回写或 \code{fsync}\\构造 request};
  \node[os small block, minimum width=24mm] (t3) at (100mm,11mm) {设备\\completion};
  \node[os small block, minimum width=26mm] (t4) at (134mm,11mm) {日志 commit\\目录同步};
  \draw[os guide] (t0.south) -- (t0.south |- 0,0);
  \draw[os guide] (t1.south) -- (t1.south |- 0,0);
  \draw[os guide] (t2.south) -- (t2.south |- 0,0);
  \draw[os guide] (t3.south) -- (t3.south |- 0,0);
  \draw[os guide] (t4.south) -- (t4.south |- 0,0);
  \node[os control, minimum width=30mm] (c1) at (0,-13mm) {完成点 1\\系统调用完成};
  \node[os control, minimum width=30mm] (c2) at (100mm,-13mm) {完成点 2\\设备请求完成};
  \node[os control, minimum width=30mm] (c3) at (134mm,-13mm) {完成点 3\\持久化发布完成};
  \draw[os guide] (c1.north) -- (c1.north |- 0,0);
  \draw[os guide] (c2.north) -- (c2.north |- 0,0);
  \draw[os guide] (c3.north) -- (c3.north |- 0,0);
  \draw[os feedback] (4mm,-4.5mm) -- node[os label, above]{掉电窗口：dirty 数据丢失} (64mm,-4.5mm);
  \draw[os feedback] (104mm,-4.5mm) -- node[os label, above]{掉电窗口：介质未必有} (131mm,-4.5mm);
\end{pdfdiagram}
\diagramnote{三个完成点在同一条真实时间线上，但彼此之间隔着可观的距离。两个红色区间是掉电会造成不同后果的窗口：第一段里数据只在内存，第二段里设备自己可能还攥着没 flush 的缓存。}

这段距离不是修辞，可以直接实测。给 64KiB 的 \code{write} 和紧随其后的 \code{fsync} 分别计时，本书实验机（aarch64 容器、闪存介质，数字只用来说明量级）上的真实输出是：

\begin{lstlisting}
write=83us fsync=3180us (wrote 65536 bytes)
write=43us fsync=195us  (wrote 65536 bytes)
write=26us fsync=140us  (wrote 65536 bytes)
\end{lstlisting}

\code{write} 几十微秒就返回，因为它只是把数据复制进 page cache 并标上 dirty；\code{fsync} 必须等设备真正确认，第一次甚至花了 3 毫秒多——完成点 1 和完成点 3 之间隔着一到两个数量级。中间那段距离不是空的，里面填满了 page cache、块层队列、驱动和设备固件各自的状态账本；本章后半部分的每一条路径，都是在解释这段距离里发生了什么。

按时间顺序展开后，三个完成点的差异更明显：

\begin{lstlisting}
write(fd, buf, len)
  -> copy from user
  -> mark page cache dirty
  -> return len

background writeback or fsync(fd)
  -> build block request
  -> map DMA buffer
  -> device completion
  -> record writeback error if any

rename("data.tmp", "data")
fsync(dirfd)
  -> publish the name
  -> make the directory update recoverable
\end{lstlisting}

三段边界各自承担不同含义。

\code{write} 返回，只说明内核按照 \code{write} 的接口语义接受了数据，或者接受了一部分数据。对普通文件来说，它经常停在 page cache 层：内存里有新内容，磁盘上未必有。

\code{fsync(fd)} 把语义往前推一层：内核要把这个文件相关的 dirty 数据和必要元数据推进到设备，并处理过程中发现的 I/O error。它比 \code{write} 强，但它通常只针对这个文件对象。

时间线里出现的 DMA 和 completion 是设备 I/O 的两个关键边界。DMA 是 direct memory access，意思是设备可以在内核授权后直接读写某段内存，不必让 CPU 一个字节一个字节搬。completion 是完成记录，表示某个提交给设备的请求结束了，内核可以根据它更新请求状态、释放 buffer、唤醒等待者。

\code{rename("data.tmp", "data")} 是发布新名字：让目录树里正式名字指向新文件。目录项本身也是元数据，所以很多崩溃安全写法还要打开目录并 \code{fsync(dirfd)}，让“名字更新”也跨过持久化边界。

如果崩溃发生在不同位置，恢复结果不同。

\begin{longtable}{@{}L{0.2\linewidth}L{0.36\linewidth}L{0.36\linewidth}@{}}
\toprule
崩溃位置 & 可能看到的状态 & 为什么不能简单说“保存成功” \\
\midrule
\code{write} 返回后、\code{fsync} 前 & 内存里有 dirty page，磁盘旧内容还在 & \code{write} 的契约不是持久化契约 \\
\code{fsync(fd)} 中途 & 部分块已写，错误可能稍后报告 & I/O error 需要被记录并返回给调用者 \\
\code{rename} 后、目录未 \code{fsync} & 文件内容可能在，但名字更新未必可恢复 & 目录项本身也是元数据，也需要提交边界 \\
日志 commit 前 & 日志里有准备信息，但不能当成功事务 & 恢复时必须丢弃未 commit 的更新 \\
日志 commit 后 & 恢复程序能重放或完成更新 & 成功来自可验证记录，不来自函数执行到最后 \\
\bottomrule
\end{longtable}

因此，“保存成功”不能只看一行返回值。它要说明数据进入了哪个对象、越过了哪个边界、失败时谁负责报告、崩溃后用什么证据恢复。page cache、块层、驱动、日志和观测工具，本质上都在追问这条时间线的不同边界。

前面的修正可以收束成一组操作系统契约：谁拥有状态，谁能修改状态，不能继续时睡在哪里，
失败如何报告，以及外部依据什么证据判断操作已经完成。

继续沿保存请求走下去。最初的函数只有局部变量和设备寄存器；一旦要支持多个程序、多个请求和失败恢复，内核必须把隐含状态变成显式对象。

\begin{longtable}{@{}L{0.2\linewidth}L{0.32\linewidth}L{0.4\linewidth}@{}}
\toprule
原来藏在哪里 & 内核对象 & 这个对象必须回答的问题 \\
\midrule
当前正在执行的函数 & task & 它能否运行，睡在哪里，怎样从同一位置恢复 \\
用户传入的指针 & VMA、PTE、page & 这个地址是否合法，读写执行权限是什么，fault 能否修复 \\
一个整数句柄 & fd table、打开文件对象 & 这个 fd 指向谁，是否可写，关闭或 exec 时怎样处理 \\
设备内部队列 & request、descriptor、completion & 哪个请求完成了，buffer 归谁，错误怎样返回 \\
缓存中的文件页 & page cache、inode & 哪些字节已修改，何时写回，谁能看到 \\
落盘顺序 & journal、commit record & 崩溃后哪些更新可信，哪些必须丢弃或重放 \\
\bottomrule
\end{longtable}

OS 对象把函数调用期间无法完整承载的状态保存到内核账本中。账本建立后，还要配套权限、
生命周期、等待队列和完成证据；缺少这些规则，系统便无法在不可信程序、慢设备和故障路径中
保持一致。

本章使用的 task、wait queue、VMA 等名称都对应具体的运行时对象。先给出对象与职责的
对照，后续章节再展开字段和状态转换：

\begin{longtable}{@{}L{0.25\linewidth}L{0.67\linewidth}@{}}
\toprule
本章叫法 & 运行时对象与职责 \\
\midrule
task & 任务对象保存执行现场、调度状态以及指向地址空间和文件表的引用 \\
wait queue & 等待队列头与等待节点保存“谁在等待哪个条件” \\
VMA & 虚拟内存区域记录一段地址范围的用途与访问策略 \\
PTE & 页表项记录当前映射、权限和处理器可见状态 \\
fd table & 文件描述符表保存整数句柄到打开文件对象的映射 \\
打开文件对象 & 保存当前偏移、打开标志、引用计数和操作方法 \\
page cache & 文件映射对象与缓存页共同保存文件内容的内存副本 \\
request & 块请求对象保存设备操作、范围、缓冲区和完成身份 \\
completion & 完成对象保存异步操作的结果并连接等待者 \\
日志提交记录 & 事务对象与提交记录界定崩溃后可以重放的更新集合 \\
\bottomrule
\end{longtable}

这里列出对象不是为了记忆结构体名称，而是为了确认每一项状态由谁负责。后面任何一章提到
“打开文件对象”时，都应当能够继续说明它保存什么状态、由谁引用以及何时释放。

\section{保存请求的状态账本}

天真的保存函数只看见一段顺序代码。真实保存请求从进入内核到对外宣布完成，中间每一步都要有对象承载状态、有边界说明语义、有证据报告失败。把这些对象排成一张账本，保存路径就不再是一串松散名词。

先把同一个保存请求放到真实 OS 语义里。用户程序看到的是几行 API：

\begin{lstlisting}
int fd = open("data.tmp", O_WRONLY | O_CREAT | O_TRUNC, 0644);
write(fd, buf, len);
fsync(fd);
close(fd);
rename("data.tmp", "data");
int dirfd = open(".", O_RDONLY | O_DIRECTORY);
fsync(dirfd);
close(dirfd);
\end{lstlisting}

这几行代码不是只能靠想象理解。可以用系统调用跟踪工具观察它大致穿过了哪些内核入口。下面的输出是删去无关参数后的示意，真实机器上路径名、fd 号码和标志细节可能不同，但边界顺序是要看的重点：

\begin{lstlisting}
$ strace -e trace=openat,write,fsync,close,renameat ./save-demo
openat(AT_FDCWD, "data.tmp", O_WRONLY|O_CREAT|O_TRUNC, 0644) = 3
write(3, "hello\n", 6)                                      = 6
fsync(3)                                                    = 0
close(3)                                                    = 0
renameat(AT_FDCWD, "data.tmp", AT_FDCWD, "data")            = 0
openat(AT_FDCWD, ".", O_RDONLY|O_DIRECTORY)                 = 3
fsync(3)                                                    = 0
close(3)                                                    = 0
\end{lstlisting}

\code{strace} 只能看到用户态和系统调用边界，不能直接告诉你 page cache、DMA 或文件系统日志里发生了什么。但它给了第一层证据：保存不是一个黑盒函数，而是一串明确的内核入口。每一行返回 \code{0} 或正数，只说明对应系统调用自己的契约成立；它不能自动替后面的边界背书。

\begin{longtable}{@{}L{0.18\linewidth}L{0.34\linewidth}L{0.4\linewidth}@{}}
\toprule
观察到的入口 & 这行真正越过的边界 & 还没有自动保证什么 \\
\midrule
\code{openat} & 路径解析、权限检查、fd 分配成功 & 后续写入一定成功 \\
\code{write} & 内核接受了 6 个字节，并推进到目标对象的写路径 & 数据已经在断电后可恢复 \\
\code{fsync(fd)} & 文件数据和必要文件元数据被要求推进到持久化边界 & 目录项名字更新已经可恢复 \\
\code{close} & fd 槽位释放，并可能暴露延迟错误 & 另一个 fd 不再引用同一个打开文件对象 \\
\code{renameat} & 目录项从临时名字发布到正式名字 & 崩溃后一定能看到新名字 \\
\code{fsync(dirfd)} & 目录对象的更新被要求推进到持久化边界 & 更上层业务事务已经提交 \\
\bottomrule
\end{longtable}

这类观察方式很重要：先看用户能写出的代码和能抓到的输出，再解释输出背后有哪些看不见的状态对象。没有这一步，\code{page cache}、\code{fd table}、\code{journal} 很容易变成互不相干的名词；有了这一步，它们都在回答同一个问题：这行返回值背后到底有什么证据。

这段代码只保留语义主线；真实工程还要检查返回值，循环处理 partial write、\code{EINTR} 和清理路径。即使只看这条主线，一个看似简单的“保存”也已经穿过了主要 OS 对象。

\begin{enumerate}
\tightlist
\item \code{open} 先走路径解析、权限检查和 fd 分配。成功后，用户拿到的只是当前进程 fd table 里的一个整数索引。
\item \code{write} 通过系统调用进入内核，复制或固定用户 buffer，把数据推进打开文件对象、page cache、socket buffer 或设备对象。返回值不等于持久化成功。
\item 如果用户页或文件页不在内存，路径可能触发 page fault、page cache miss、块 I/O 和调度睡眠。
\item 若块设备异步写入，内核要建立 request，映射 DMA buffer，提交描述符，等待 completion，再解除映射。
\item \code{fsync} 把“内核已经接受写入”推进到更强的持久化边界，但仍要处理设备错误、flush、日志提交和元数据顺序。
\item \code{rename} 把临时文件发布成正式名字；目录 \code{fsync} 让名字更新本身也有崩溃恢复语义。
\end{enumerate}

把这些状态写成一张账本：

\begin{longtable}{@{}L{0.16\linewidth}L{0.24\linewidth}L{0.28\linewidth}L{0.24\linewidth}@{}}
\toprule
阶段 & 状态对象 & 关键边界 & 失败证据 \\
\midrule
入口 & task、入口现场 & 用户态切到内核态，保存返回位置 & 错误码、审计记录 \\
授权 & fd、打开文件对象、inode & fd 是否存在，目标是否可写 & \code{-EBADF}、\code{-EACCES} \\
复制 & VMA、PTE、page & 用户 buffer 是否可读 & \code{-EFAULT}、部分写入 \\
缓存 & page cache、inode & dirty 页只是内存状态 & dirty/writeback 计数、回写错误 \\
提交 I/O & request、DMA map & buffer 在 completion 前不能释放 & timeout、I/O error、completion status \\
持久化 & journal、flush、directory entry & 数据和名字是否崩溃后仍可解释 & fsync 错误、replay、fsck 日志 \\
\bottomrule
\end{longtable}

这张表把保存请求从一个 API 调用拆成了六个阶段：入口、授权、复制、缓存、异步提交和持久化发布。任何一个阶段说不清，程序就可能在错误路径里撒谎：返回了成功，但对象状态并没有达到用户以为的边界。

\begin{warningbox}
保存路径里最容易混淆的五个边界：
\begin{itemize}
\tightlist
\item \code{write} 返回成功，不等于数据已经落盘。
\item 线程处于 runnable，不等于正在运行。
\item fd 是整数句柄，不等于文件本体。
\item 虚拟地址落在 VMA 内，不等于 PTE 已经存在。
\item 唤醒函数被调用，不等于等待者已经继续执行。
\end{itemize}
\end{warningbox}

再把这段保存代码当成一个完整案例看。很多程序会先写成这样：

\begin{lstlisting}
bool save_bad(const char* path, const char* buf, size_t len) {
  int fd = open(path, O_WRONLY | O_CREAT | O_TRUNC, 0644);
  if (fd < 0) return false;
  if (write(fd, buf, len) != len) return false;
  close(fd);
  return true;
}
\end{lstlisting}

它的表面意思很直接：打开目标文件，截断旧内容，写入新内容，关闭，返回成功。但把前面的账本套上去，会发现它在三个位置同时冒险。

\begin{longtable}{@{}L{0.2\linewidth}L{0.36\linewidth}L{0.36\linewidth}@{}}
\toprule
位置 & 发生了什么 & 失败后用户可能看到什么 \\
\midrule
\code{O\_TRUNC} & 旧文件长度先变成 0 & 后续写失败时，旧内容已经丢失 \\
一次 \code{write} & 可能只复制部分字节，或停在 page cache & 返回短写时，新文件是半成品 \\
\code{close} 后返回 & 未检查 close/fsync 的延迟错误 & 后台写回失败，用户却看到保存成功 \\
\bottomrule
\end{longtable}

把它放到一个具体文件上看，会更容易发现危险。假设正式文件 \code{data} 原来保存的是 \verb|"old\n"|，这次要保存的新内容是 \verb|"new\n"|。用户想要的结果其实很严格：失败时最好还能读到旧版本；成功时读到完整新版本；最糟糕的状态是旧版本被破坏，而新版本又不完整。

\begin{longtable}{@{}L{0.22\linewidth}L{0.34\linewidth}L{0.36\linewidth}@{}}
\toprule
执行到哪里 & \code{data} 可能处于什么状态 & 为什么危险 \\
\midrule
\code{open(... O\_TRUNC)} 之后 & 长度已经变成 0 & 还没写任何新数据，旧版本已经被破坏 \\
\code{write} 短写之后 & 可能只有 \code{"ne"} 这样的前缀 & 系统调用允许部分推进，不能假设一次写完 \\
\code{write} 返回、未 \code{fsync} & page cache 里可能是新内容，磁盘上可能仍是旧内容或部分更新 & 函数返回和持久化不是同一个边界 \\
\code{close} 未检查错误 & fd 被释放，但延迟 I/O 错误可能被忽略 & 调用者看到 true，实际保存可能失败 \\
\bottomrule
\end{longtable}

改进方向不是机械记住“要用临时文件”，而是从失败路径推出临时文件的必要性：旧文件不能在新内容可靠之前被破坏。更稳的写法先把新内容写进同目录下的临时文件，确认它的数据和必要元数据已经推进，再用 \code{rename} 一次性替换名字。

先把这个推导拆开。错误版本从正式文件本身下手，所以第一步 \code{O\_TRUNC} 就已经破坏旧版本；后面任何失败都会让用户处在尴尬状态：旧内容没了，新内容也没完整提交。临时文件版本把“准备新内容”和“发布新名字”分成两个阶段。准备阶段只影响临时名字，失败就删除临时文件；发布阶段用 \code{rename} 切换目录项，使正式名字尽量只在旧版本和新版本之间切换。

\begin{longtable}{@{}L{0.2\linewidth}L{0.35\linewidth}L{0.37\linewidth}@{}}
\toprule
设计选择 & 它来自哪个失败 & 边界含义 \\
\midrule
不用 \code{O\_TRUNC} 直接改正式文件 & 写到一半失败会破坏旧版本 & 旧文件在新版本准备好前保持可用 \\
使用同目录临时文件 & 跨文件系统 \code{rename} 不能保证同样语义 & 临时文件和正式文件处在同一个目录事务范围内 \\
循环 \code{write\_all} & 一次 \code{write} 可能短写或被信号打断 & 未写完整就不能进入发布阶段 \\
先 \code{fsync(tmpfd)} & page cache 中的新内容不等于持久内容 & 新文件内容先越过文件持久化边界 \\
再 \code{rename} & 正式名字只能指向旧版本或新版本 & 发布动作集中到目录项切换 \\
最后 \code{fsync(dirfd)} & 目录项本身也是元数据 & 崩溃后正式名字更新有恢复证据 \\
\bottomrule
\end{longtable}

先把 \code{write\_all} 摊开。它不是为了把代码写得“保险一点”，而是因为 \code{write} 的基本契约就是“返回本次实际推进了多少字节”。只要返回值小于请求长度，调用者就不能假装完整数据已经进入保存路径。

\begin{lstlisting}
bool write_all(int fd, const char* buf, size_t len) {
  size_t off = 0;

  while (off < len) {
    ssize_t n = write(fd, buf + off, len - off);

    if (n > 0) {
      off += (size_t)n;
      continue;
    }

    if (n == 0) {
      return false;       // no progress
    }

    if (errno == EINTR) {
      continue;           // interrupted before a final result
    }

    if (errno == EAGAIN || errno == EWOULDBLOCK) {
      wait_until_writable(fd);
      continue;
    }

    return false;         // ENOSPC, EIO, EFAULT, ...
  }

  return true;
}
\end{lstlisting}

假设要写入 6 个字节 \verb|hello\n|。第一次 \code{write} 只返回 2，表示 \code{he} 已经推进；第二次被信号打断，返回 \code{-1/EINTR}，表示没有得到一个最终完成结果；第三次返回 4，才把剩下的 \verb|llo\n| 推进。此时 \code{off} 的变化是保存程序必须记住的状态。

\begin{longtable}{@{}L{0.18\linewidth}L{0.28\linewidth}L{0.22\linewidth}L{0.24\linewidth}@{}}
\toprule
调用 & 返回 & \code{off} 变化 & 保存程序该做什么 \\
\midrule
第 1 次 \code{write} & 2 & 0 变成 2 & 继续写剩余 4 字节 \\
第 2 次 \code{write} & \code{-1/EINTR} & 仍是 2 & 不能报成功，重试同一段剩余数据 \\
第 3 次 \code{write} & 4 & 2 变成 6 & 才能进入 \code{fsync(tmpfd)} \\
\bottomrule
\end{longtable}

\code{EINTR}、\code{EAGAIN} 和短写看起来像细枝末节，但它们对应三种不同边界。\code{EINTR} 表示调用被信号路径打断；\code{EAGAIN} 表示对象现在不能推进，调用者要等待可写条件；短写表示系统已经接受了一部分数据，下一次必须从新 offset 继续。若保存程序不记录 \code{off}，就会重复写前缀、漏写后缀，或者把半成品当完整版本发布。

\begin{lstlisting}
bool save_better(const char* dir, const char* name,
                 const char* buf, size_t len) {
  int dirfd = open(dir, O_RDONLY | O_DIRECTORY);
  if (dirfd < 0) return false;

  char tmp_name[64];
  int tmpfd = open_temp_at(dirfd, tmp_name, sizeof(tmp_name));
  if (tmpfd < 0) goto fail_dir;

  if (!write_all(tmpfd, buf, len)) goto fail_tmp;
  if (fsync(tmpfd) < 0) goto fail_tmp;
  if (close(tmpfd) < 0) {
    tmpfd = -1;
    goto fail_unlink;
  }
  tmpfd = -1;

  if (renameat(dirfd, tmp_name, dirfd, name) < 0) goto fail_unlink;
  if (fsync(dirfd) < 0) goto fail_dir;

  close(dirfd);
  return true;

fail_tmp:
  close_if_open(tmpfd);
fail_unlink:
  unlinkat(dirfd, tmp_name, 0);
fail_dir:
  close(dirfd);
  return false;
}
\end{lstlisting}

这段代码不是完整库函数；真实工程还要处理 \code{EINTR}、权限、临时文件命名、保留文件模式、跨文件系统 rename 等细节。但它已经把关键边界摆清楚了。

逐行看它的状态推进会更清楚。打开目录不是多余动作，因为最后要同步目录项；创建临时文件不是为了名字好看，而是为了让新内容有独立承载对象；\code{write\_all} 不是普通循环，而是在处理“接口允许短推进”这个契约；\code{fsync(tmpfd)} 把临时文件内容向持久化推进；\code{close(tmpfd)} 还可能暴露延迟错误；\code{renameat} 才把正式名字切到新文件；\code{fsync(dirfd)} 则让这个名字变化本身在崩溃后可恢复。

\begin{longtable}{@{}L{0.19\linewidth}L{0.35\linewidth}L{0.38\linewidth}@{}}
\toprule
处理阶段 & 此刻系统必须记住什么 & 如果失败该怎么解释 \\
\midrule
\code{open(dir)} & 目录对象和权限 & 目录打不开，保存根本没有进入准备阶段 \\
\code{open\_temp\_at} & 临时 inode、临时目录项、fd 引用 & 只清理临时对象，不影响正式文件 \\
\code{write\_all} & 已写入字节数、dirty page、可能的短写 & 未完整写入，不允许发布 \\
\code{fsync(tmpfd)} & 数据和必要元数据的提交结果 & 返回错误说明新版本不可信 \\
\code{renameat} & 目录项从旧 inode 切到新 inode & 失败时正式名字仍应指向旧版本 \\
\code{fsync(dirfd)} & 目录项更新的持久化证据 & 失败时必须向上报告，不能宣称完全保存 \\
\bottomrule
\end{longtable}

\begin{longtable}{@{}L{0.18\linewidth}L{0.38\linewidth}L{0.36\linewidth}@{}}
\toprule
动作 & 它保护什么 & 崩溃后恢复含义 \\
\midrule
同目录临时文件 & 保证 rename 不跨文件系统，旧名字暂时不变 & 旧文件仍可作为最后成功版本 \\
\code{write\_all} & 处理短写、\code{EINTR}、\code{EAGAIN} 等推进问题 & 临时文件要么完整写入，要么不发布 \\
\code{fsync(tmpfd)} & 推进临时文件数据和必要元数据 & 新内容有机会在崩溃后被读回 \\
\code{rename} & 原子切换目录项指向 & 名字层面看到旧版本或新版本，不应看到半个名字 \\
\code{fsync(dirfd)} & 推进目录项更新本身 & 崩溃后正式名字更新有恢复证据 \\
\bottomrule
\end{longtable}

再用同一个 \verb|data = "old\n"|、新内容 \verb|"new\n"| 的例子看崩溃点。下面的表不假设某一种具体文件系统实现细节，只抓通用语义边界：正式名字什么时候还能指向旧版本，临时文件什么时候可以被丢弃，什么时候才能对用户说“新版本已经发布”。

\begin{longtable}{@{}L{0.2\linewidth}L{0.34\linewidth}L{0.38\linewidth}@{}}
\toprule
崩溃发生在 & 重启后合理看到的状态 & 保存程序应该如何解释 \\
\midrule
临时文件创建前 & \code{data} 仍是旧版本 & 保存没有开始影响正式文件 \\
\code{write\_all} 中途 & \code{data} 仍是旧版本，可能残留半个临时文件 & 临时文件不能发布，清理后重试 \\
\code{fsync(tmpfd)} 失败 & \code{data} 仍是旧版本，新内容不可信 & 必须返回失败，不能继续 rename \\
\code{renameat} 前 & \code{data} 仍指向旧版本，临时文件可能完整 & 新版本还没有发布给正式名字 \\
\code{renameat} 后、目录未同步 & 运行时能看到新名字，但崩溃恢复后名字更新可能没有证据 & 不能把目录项持久化当成已经完成 \\
\code{fsync(dirfd)} 成功后 & \code{data} 应当指向可恢复的新版本 & 可以向上层报告保存跨过发布边界 \\
\bottomrule
\end{longtable}

这张表比“用临时文件更安全”更重要。它说明安全不是来自某个函数名，而是来自一组顺序关系：先保护旧版本，再准备新版本，再持久化新版本，再发布名字，最后持久化名字。任何一步失败，都要停在还能解释的状态上。

这段保存代码把前面的对象串成了一条完整链路。task 解释了为什么等待 I/O 时不能忙等；VMA/PTE 解释了 \code{buf} 为什么要复制检查；fd、打开文件对象和 inode 解释了为什么整数句柄不是文件本体；page cache 解释了 \code{write} 为什么可以早于设备完成返回；request/DMA/completion 解释了内核如何跟踪异步设备结果；fsync/rename/目录 fsync 解释了“保存成功”如何从内存状态变成崩溃后可解释的证据。

再看“对象”这个词，它就不再是抽象名词。对象就是内核为了记住某类状态而建立的账本。前面保存案例已经见过几本账：

\begin{longtable}{@{}L{0.18\linewidth}L{0.32\linewidth}L{0.42\linewidth}@{}}
\toprule
对象 & 它记住什么 & 没有它会怎样 \\
\midrule
task & 当前执行流是否 running、runnable、blocked，睡在哪里 & 设备慢时只能忙等，或者睡了以后无法恢复 \\
wait queue & 哪些 task 在等同一个谓词 & 设备完成后不知道该叫醒谁，容易 lost wakeup \\
VMA/PTE & 地址区间的意图和单页翻译事实 & 用户随便传一个指针，内核不知道该复制、补页还是拒绝 \\
fd table/打开文件对象 & 整数 fd 指向哪个打开对象，offset 和读写权限是什么 & \code{dup}、\code{close}、共享 offset 都解释不清 \\
page cache & 文件内容在内存中的页面，哪些是 dirty & \code{write} 返回后无法说明数据停在哪里 \\
请求/完成记录 & 哪次设备 I/O 已提交，哪次已经完成 & 中断来了也不知道对应哪个保存请求 \\
日志/提交记录 & 哪些更新崩溃后可信 & 写一半掉电后无法区分成功更新和半成品 \\
\bottomrule
\end{longtable}

这张表背后的判断规则是：只要某个状态会被多个路径共享，或者失败后需要恢复，就不能只藏在一个函数局部变量里。内核必须给它名字、所有者、修改入口、等待位置和失败证据。

同样的账本形状也会出现在别的对象上。socket 会把网络连接变成对象；地址空间会把整套页表和 VMA 组织起来；块层会把一次大写入拆成多个 request；安全子系统会把权限判断变成可审计的拒绝路径。对象一旦进入内核，就必须说明它记录什么、谁能改它、错了以后怎么证明。

保存案例还暴露出一个分层事实：同一个词在用户 API、内核对象和硬件事件三层里含义不同。用户说地址，可能只是一个指针值；内核说地址，要看 VMA、PTE 和权限；硬件说地址，要看页表翻译、TLB 和物理地址。用户说完成，可能只是函数返回；内核说完成，可能只是 page cache 或 request 状态更新；设备说完成，可能只是某次 I/O 结束。

因此，同一个动作要同时放进三层：用户 API 层、内核对象层、硬件事件层。

\begin{longtable}{@{}L{0.17\linewidth}L{0.27\linewidth}L{0.28\linewidth}L{0.2\linewidth}@{}}
\toprule
说法 & 用户 API 层 & 内核对象层 & 硬件层 \\
\midrule
地址 & 指针值，例如 \code{0x1000} & VMA/PTE 和权限 & 页表翻译、物理地址、设备寄存器地址 \\
可读 & \code{read} 不会立刻阻塞 & buffer/队列中有数据或条件满足 & 中断、completion、DMA 写入 \\
写完成 & 系统调用返回字节数 & page cache 或 request 状态更新 & 设备接受请求，介质 flush 可能还没发生 \\
上下文 & 用户调用链和线程身份 & task、保存的返回位置、内核栈 & 程序计数器、栈指针、寄存器 \\
错误 & errno 或信号 & 对象状态拒绝、权限失败、生命周期结束 & fault、设备错误、中断错误 \\
\bottomrule
\end{longtable}

以“写完成”为例。写普通文件时，用户看到 \code{write} 返回；内核看到 page cache 被改成 dirty；设备可能还什么都没看到。等后台回写或 \code{fsync} 提交 request 后，设备才开始处理这次 I/O。等 completion 到来，内核知道这次设备请求结束了。等 flush 或日志 commit 过了，崩溃恢复才有证据。这些都叫“完成”，但边界完全不同。

再看“可读”。普通文件几乎总是可读，因为内核可以从 page cache 或磁盘推进；socket 可读表示接收队列有数据、连接关闭或错误；设备 fd 可读可能表示驱动队列里有 completion。用户只看到 \code{poll} 返回，但内核必须把对象状态和硬件事件翻译成统一接口。

\section{从硬件事件到内核决策}

一次保存请求不是一条顺着函数调用栈走到底的直线。它会被四类入口打断和续上：系统调用入口、缺页入口、定时器入口和设备完成入口。每条路径都从硬件事实进入内核决策，再回到用户能观察到的结果。

\begin{pdfdiagram}
  \node[os memory, minimum width=36mm] (h1) at (0,0) {syscall 指令\\特权级切换};
  \node[os memory, minimum width=36mm] (h2) at (0,-14mm) {页表检查失败\\保存 fault 地址};
  \node[os memory, minimum width=36mm] (h3) at (0,-28mm) {定时器触发\\指令边界采样};
  \node[os memory, minimum width=36mm] (h4) at (0,-42mm) {completion\\设备中断};
  \node[os state, minimum width=44mm, minimum height=56mm] (k) at (64mm,-21mm) {};
  \node[os title, align=center] at (64mm,-21mm) {内核决策\\同一套对象账本\\task、VMA、fd\\request、wait queue};
  \node[os compute, minimum width=38mm] (r1) at (126mm,0) {返回字节数\\或 \code{-errno}};
  \node[os compute, minimum width=38mm] (r2) at (126mm,-14mm) {重试指令\\或发送信号};
  \node[os compute, minimum width=38mm] (r3) at (126mm,-28mm) {选择下一个\\task 运行};
  \node[os compute, minimum width=38mm] (r4) at (126mm,-42mm) {更新 request\\唤醒等待者};
  \draw[os strong bus] (h1.east) -- (k.west |- h1.east);
  \draw[os strong bus] (h2.east) -- (k.west |- h2.east);
  \draw[os strong bus] (h3.east) -- (k.west |- h3.east);
  \draw[os strong bus] (h4.east) -- (k.west |- h4.east);
  \draw[os strong bus] (k.east |- r1.west) -- (r1.west);
  \draw[os strong bus] (k.east |- r2.west) -- (r2.west);
  \draw[os strong bus] (k.east |- r3.west) -- (r3.west);
  \draw[os strong bus] (k.east |- r4.west) -- (r4.west);
\end{pdfdiagram}
\diagramnote{四条路径入口不同、结果不同，但中间经过的是同一套内核对象账本。理解其中任何一条，都要能说清它读写了哪些共享状态——这正是后面每一章的展开方式。}

第一条路径是系统调用入口。应用看见的是一个 C 库函数，内核看见的是一次受控入口和多个对象状态更新。这里的寄存器名字来自 x86-64 调用约定，用来表达四件事：系统调用号、fd、buffer 地址和长度被交给内核。

\begin{lstlisting}
user:
  rax = SYS_write
  rdi = fd
  rsi = user_buf
  rdx = len
  syscall

kernel entry:
  save user return position and registers
  switch to kernel-owned stack
  locate current task and fd table
  validate fd refers to writable file, socket, or pipe
  validate user buffer range and copy or pin pages
  update target object state
  return byte count or -errno in rax
\end{lstlisting}

\begin{longtable}{p{0.2\linewidth}p{0.36\linewidth}p{0.32\linewidth}}
\toprule
阶段 & 关键检查 & 失败表达 \\
\midrule
系统调用入口 & syscall number 是否存在，当前上下文是否允许 & \code{-ENOSYS}、拒绝或信号 \\
fd 查找 & fd 是否打开，是否有写权限 & \code{-EBADF}、\code{-EPERM} \\
用户 buffer & 地址格式是否有效、VMA 是否允许读、复制中是否 fault & \code{-EFAULT} 或部分写入 \\
目标对象 & 文件、pipe 或 socket 是否可写、是否阻塞 & \code{-EAGAIN}、睡眠、信号中断 \\
提交语义 & 数据进入哪里：page cache、socket buffer、设备队列 & 返回值不一定表示落盘或发出网络包 \\
\bottomrule
\end{longtable}

这条路径说明了 OS 契约的层次。系统调用返回成功，只说明内核接受并完成了该接口定义下的一部分工作。写普通文件可能只是写入 page cache；写 socket 可能只是进入发送缓冲；写终端可能经过驱动队列和中断。若业务需要“崩溃后仍存在”，还要继续追 \code{fsync} 和存储 flush；若业务需要“对端已收到”，还要追协议确认。

第二条路径是 page fault。内核按用户地址复制 \code{user\_buf} 时，CPU 可能发现页表里没有可用翻译，于是触发 page fault。硬件做的是检查和转交，内核做的是解释和修复。

\begin{lstlisting}
CPU:
  compute the address being loaded
  check whether the address form is valid
  lookup TLB or walk page table
  if missing or permission denied:
    save faulting instruction position and error code
    record fault address
    enter page fault handler

kernel:
  find VMA covering fault address
  compare access type with VMA permissions
  inspect PTE state
  allocate a page, read a file page, copy a shared page, or reject
  install PTE and clear stale address-cache entries when needed
  return to the faulting instruction or deliver signal
\end{lstlisting}

\begin{longtable}{p{0.22\linewidth}p{0.34\linewidth}p{0.32\linewidth}}
\toprule
硬件看到 & 内核进一步区分 & 结果 \\
\midrule
present=0 & 地址属于匿名 VMA & 分配零页，minor fault \\
present=0 & 地址属于文件映射 & page cache 命中或发起磁盘 I/O \\
write denied & PTE 标记 shared-copy & 复制页并恢复写权限 \\
execute denied & NX 或 VMA 不可执行 & 发送信号，安全拒绝 \\
user bit denied & 用户访问内核专用页 & 发送信号或记录攻击/bug \\
\bottomrule
\end{longtable}

缺页机制的关键不是“慢”，而是“可恢复边界”。如果 CPU 不能保存出错指令的位置，修好页表后就无法从原指令重试；如果内核不能判断 VMA 意图，present=0 就无法区分合法懒分配和非法访问；如果旧的地址翻译缓存没有清掉，PTE 更新后旧权限仍可能生效。

第三条路径是定时器中断。用户程序可能永远不主动让出 CPU。定时器中断给内核一个周期性入口，让它能统计时间、处理超时并决定是否切换任务。

\begin{lstlisting}
hardware timer fires:
  CPU takes interrupt at instruction boundary
  entry saves interrupted return position and registers
  timer handler increments scheduler clock
  expired timers wake sleeping tasks
  current task runtime is charged
  if time slice or policy says preempt:
    set need_resched
  return path calls scheduler before going back to user
  scheduler saves current kernel context and restores next
\end{lstlisting}

这里有两个容易混淆的点。第一，中断处理函数通常不在任意深处直接把用户栈改成另一个任务；它设置调度标志，在安全返回点进入调度。第二，切换任务不是只改一个 current 指针，还可能涉及内核栈、需要保留的寄存器、地址空间和地址翻译缓存。

\begin{longtable}{p{0.22\linewidth}p{0.34\linewidth}p{0.32\linewidth}}
\toprule
动作 & 需要保存什么 & 保存错误的现象 \\
\midrule
中断进入 & 用户返回位置、栈指针和寄存器 & 返回后用户程序寄存器随机坏 \\
调度统计 & 当前任务运行时间和状态 & CPU 时间不公平或饥饿 \\
切换内核上下文 & 内核栈指针、必须保留的寄存器 & 任务从错误调用链恢复 \\
切换地址空间 & 当前页表身份、必要的翻译缓存处理 & 任务访问到别人的内存 \\
\bottomrule
\end{longtable}

第四条路径是设备完成。设备 I/O 是异步的。提交请求后，任务可能睡眠；设备通过 DMA 和中断报告完成；内核再把 completion 变成任务可见结果。

\begin{lstlisting}
submit I/O:
  allocate request object
  map buffer for DMA
  fill device-visible descriptor
  publish descriptor in the required order
  notify the device
  if synchronous API:
    put task on wait queue and schedule away

device:
  reads descriptor by DMA
  transfers data
  writes completion entry
  raises interrupt

interrupt/bottom half:
  read completion queue
  unmap DMA and mark request done
  store status and byte count
  wake tasks waiting on request
\end{lstlisting}

这条路径解释了为什么 I/O 子系统一定有 request identity，也就是“请求身份”。中断到来时，CPU 不知道“哪个用户函数在等它”，只看到某个设备队列有 completion。内核必须通过队列位置、tag、request 指针或 completion entry 找回原请求，撤销 DMA 授权，更新上层对象，再唤醒等待者。没有这层身份映射，异步 I/O 无法可靠对应到发起者。

把这个过程压到一个设备队列里看。驱动把请求写进一块设备也能读取的 descriptor ring；descriptor 是一条描述符，里面记录 buffer 地址、长度、方向和 tag。doorbell register 是设备暴露给 CPU 的一个通知寄存器；驱动把 descriptor 写好后，再写 doorbell，告诉设备“队列里有新请求”。completion queue 则是设备写回结果的队列。

\begin{lstlisting}
submit A:
  descriptor[5] = { tag: 17, buffer: page0, len: 4096, op: write }
  doorbell = 5

submit B:
  descriptor[6] = { tag: 18, buffer: page1, len: 4096, op: write }
  doorbell = 6

completion queue:
  { tag: 18, status: OK,  bytes: 4096 }
  { tag: 17, status: EIO, bytes: 0    }
\end{lstlisting}

完成顺序不一定等于提交顺序。设备可以先完成 B，再报告 A 失败。若内核只按“第一个 completion 对应第一个等待者”处理，就会把 B 的成功错还给 A，把 A 的错误错还给 B。tag 的作用就是把设备事实重新接回内核对象账本。

\begin{pdfdiagram}
  \node[os compute, minimum width=26mm] (reqA) at (0,0) {request A\\tag 17};
  \node[os compute, minimum width=26mm] (reqB) at (0,-13mm) {request B\\tag 18};
  \node[os memory, minimum width=46mm] (ring) at (56mm,-6.5mm) {descriptor ring\\{[}5{]}: tag 17, write, page0\\{[}6{]}: tag 18, write, page1};
  \node[os control, minimum width=26mm] (bell) at (56mm,-21mm) {doorbell 寄存器};
  \node[os memory, minimum width=30mm] (dev) at (112mm,-6.5mm) {设备\\DMA 取描述符};
  \node[os state, minimum width=46mm] (cq) at (56mm,-38mm) {completion queue\\tag 18: OK, 4096 字节\\tag 17: EIO, 0 字节};
  \draw[os strong bus] (reqA.east) -- (ring.west);
  \draw[os strong bus] (reqB.east) -- (ring.west);
  \draw[os bus] (ring) -- (bell);
  \draw[os strong bus] (ring) -- node[os label, above]{DMA 读} (dev);
  \draw[os bus] (bell.east) -| node[os label, near start, below]{通知} (dev.south);
  \draw[os strong bus] (dev.south) |- node[os label, near end, right]{DMA 写} (cq.east);
  \draw[os feedback] (cq.west) -- (reqB.south east);
  \draw[os feedback] (cq.west) -- (reqA.south east);
  \node[os label] at (14mm,-27mm) {红色虚线：按 tag 找回原 request};
\end{pdfdiagram}
\diagramnote{完成顺序不等于提交顺序：B 先成功、A 后失败是合法结果。completion entry 里的 tag 把设备事实重新接回正确的 request 账本，两条红色虚线就是这次“找回”。}

\begin{longtable}{@{}L{0.2\linewidth}L{0.36\linewidth}L{0.36\linewidth}@{}}
\toprule
阶段 & 硬件/驱动看到什么 & 内核必须维护什么 \\
\midrule
填 descriptor & buffer 地址、长度、操作类型、tag & request 对象仍然活着，buffer 不能释放 \\
写 doorbell & 设备被通知去取 descriptor & descriptor 字段必须已经完整可见 \\
设备 DMA & 设备直接读写内存 buffer & 页面要被固定或映射，防止被回收复用 \\
写 completion & completion entry 带回 tag 和 status & 用 tag 找回 request，记录成功或错误 \\
唤醒等待者 & 等待队列里的 task 重新变 runnable & 醒来后读取 request status，不猜测结果 \\
\bottomrule
\end{longtable}

这里最危险的错误不是“设备慢”，而是生命周期错位。若 completion 之前释放 buffer，设备可能把数据写进已经分配给别人的页面；若 completion 没有 status，调用者只能知道“有中断”，不知道哪次写失败；若唤醒早于 request status 更新，等待者醒来仍可能读到旧状态。异步设备因此迫使 OS 明确三件事：请求身份、buffer 生命周期和完成证据。

四条路径讲完，可以进一步比较它们进入内核时携带的证据和最终交给哪类机制处理：

\begin{longtable}{@{}L{0.18\linewidth}L{0.38\linewidth}L{0.36\linewidth}@{}}
\toprule
路径 & 进入内核时的主要证据 & 后续处理 \\
\midrule
系统调用 & 调用号、参数寄存器、用户返回现场 & 分发、参数验证、权限检查与返回 \\
缺页异常 & 故障地址、访问类型、当前页表状态 & VMA 查找、权限判断、建立映射或报告错误 \\
定时器中断 & 当前任务、时间账本和抢占状态 & 计时、调度决策与可能的任务切换 \\
设备完成 & 完成队列项、请求标识和状态码 & 找回请求、转移缓冲区所有权并唤醒等待者 \\
\bottomrule
\end{longtable}

四条路径的共同点是：入口只提供作出决定所需的证据，真正的语义由内核对象和状态机完成。
后续章节会分别展开这些决策如何保持权限、生命周期与失败边界。

\section{从硬件约束到操作系统抽象}

内核账本必须建立在硬件提供的条件之上。CPU 能保存执行位置，因此 task 可以睡眠后恢复；页表能够检查地址权限，因此内核可以拒绝非法 buffer；设备以异步方式完成请求，因此内核需要维护 request 和 completion；存储可能在任意时刻失去供电，因此文件系统必须定义提交边界。

\begin{longtable}{p{0.22\linewidth}p{0.34\linewidth}p{0.32\linewidth}}
\toprule
硬件条件 & 操作系统抽象 & 关键问题 \\
\midrule
程序计数器、寄存器和栈可保存 & task 上下文 & 切走后如何从同一位置恢复 \\
特权级限制敏感指令 & 用户态/内核态 & 用户如何安全请求内核服务 \\
异常入口能保存出错位置 & page fault、signal、debug & 错误能否修复并返回原指令 \\
定时器可强制中断 & 抢占式调度 & 不合作任务如何被收回 CPU \\
MMU 检查页表权限 & 地址空间 & 进程如何互相隔离又能共享 \\
cache/TLB 需要维护 & 性能和一致性策略 & 为什么迁移、锁和取消映射有成本 \\
设备寄存器控制设备 & 驱动模型 & 寄存器副作用和顺序如何管理 \\
DMA 异步访问内存 & 高性能 I/O & buffer 何时归 CPU，何时归设备 \\
\bottomrule
\end{longtable}

看到“进程”，需要同时看到一组保存的 CPU 上下文、地址空间和内核对象，而不只是 pid。看到“文件”，需要同时看到 fd 表、打开文件对象、inode、page cache、块层和设备队列，而不只是路径。看到“权限”，需要同时看到 CPU 特权级、页表权限、文件权限和对象引用，而不只是口令或用户 ID。

用户接口词、内核对象词和硬件词经常指向同一条路径的不同层。用户接口会说 process、file descriptor、signal、pipe、socket；内核对象会说 task、file、inode、wait queue；硬件路径会说寄存器、页表、中断、DMA、设备寄存器。三套词不是互相替代，而是共同描述一次动作怎样穿过系统。

\begin{longtable}{@{}L{0.18\linewidth}L{0.25\linewidth}L{0.27\linewidth}L{0.22\linewidth}@{}}
\toprule
用户语义 & 内核对象 & 硬件支点 & 关键边界 \\
\midrule
process & task、地址空间、fd table & 寄存器、页表、定时器 & 身份、地址空间、权限分离 \\
thread & task、内核栈、共享地址空间 & 寄存器、栈指针 & 可调度上下文，不复制全部资源 \\
fd & fd table、打开文件对象 & 设备队列、page cache、DMA & 整数句柄不是文件本体 \\
signal & 待处理信号、用户返回现场 & 保存的异常/中断现场 & 异步交付但必须恢复用户现场 \\
socket & socket 对象、wait queue & 网卡队列、中断、DMA & 协议状态和设备完成分离 \\
timer & 定时器对象、等待队列 & 硬件定时器中断 & 未来事件需要硬件唤醒入口 \\
\bottomrule
\end{longtable}

例如 fd 是一个用户态整数，但内核要通过当前进程的 fd table 找到打开文件对象，再通过 file operations 调到 pipe、socket、inode 或 device 的具体实现。若底层是网卡，最终还会走网络缓冲区、驱动队列、DMA 和 completion。把 fd 当成“文件本体”，就解释不了 \code{dup} 后共享 offset、socket readiness、设备拔出后旧 fd 返回错误这些现象。

单片机和裸机程序能暴露 OS 的原始需求。一个裸机程序拥有全部内存和设备，所有代码同权，主循环靠自觉返回，中断直接修改全局状态。高级机制出现之前，失败通常已经在这种小系统里出现过。

\begin{longtable}{p{0.24\linewidth}p{0.32\linewidth}p{0.32\linewidth}}
\toprule
裸机阶段的问题 & 局部处理 & 通用 OS 机制 \\
\midrule
函数忙等设备 & 改成状态机或中断 & 阻塞、等待队列、completion \\
任务靠自觉返回 & 任务表和显式 yield & 线程和协作式调度 \\
任务可能不合作 & 加硬件定时器 & 抢占和时间片 \\
全局变量被中断处理函数打断 & 关中断临界区 & 锁、原子操作、内存屏障 \\
所有代码同权 & 约定模块边界 & CPU 特权级和系统调用 \\
所有内存同权 & 手工内存布局 & MMU、地址空间、页表权限 \\
CPU 搬每个字节 & DMA 控制器 & 异步 I/O 和驱动队列 \\
\bottomrule
\end{longtable}

如果不知道硬件能否强制中断、能否区分用户/内核、能否检查页表权限，就无法判断一个 OS 设计到底是“规范约定”还是“硬件保证”。

\section{保存路径的五个边界}

保存请求走到这里，已经不是一段顺序代码，而是一串状态和边界。每个边界都回答一个具体问题：对象是谁，谁拥有它，哪个入口能修改它，不能继续时停在哪里，完成时用什么证据说明结果可信。

\begin{longtable}{p{0.18\linewidth}p{0.74\linewidth}}
\toprule
问题 & 保存路径中的内容 \\
\midrule
状态对象 & 进程、线程、页、fd、inode、socket、设备请求、日志记录 \\
所有者 & 用户态对象、内核对象、硬件寄存器、设备控制器 \\
入口 & 系统调用、异常、中断、内核线程、后台回收 \\
权限 & 特权级、页表权限、文件权限、资源能力、命名空间、引用计数 \\
等待 & runqueue、futex、page fault、I/O completion、锁、条件变量 \\
提交 & 状态更新、引用释放、写入 page cache、落盘、目录项发布生效 \\
证据 & 返回值、errno、trace、计数器、日志、测试断言 \\
\bottomrule
\end{longtable}

例如一次 \code{write(fd, buf, len)} 可以这样拆：用户态只有一个整数 fd 和一段用户 buffer；内核要先从进程 fd 表找到打开文件对象，再检查写权限，再复制用户 buffer，再把页标记为 dirty，最后返回写入字节数或错误。这个返回值通常不等于数据已经安全落盘。要证明持久化，还需要 \code{fsync}、目录同步或更高层提交记录。

再看一次 page fault。用户态执行 load/store，硬件发现页表项不存在或权限不符，CPU 进入异常入口，内核根据 VMA 判断这是合法缺页、权限错误还是非法地址。合法缺页可能分配物理页或读取文件页；权限错误可能发送信号；非法访问可能终止进程。这里的核心不是“缺页很慢”，而是虚拟地址访问被拆成硬件检查和内核决策。

只看正常路径仍然不够。OS 机制真正麻烦的地方，是反例会从边界处钻出来。

\begin{longtable}{@{}L{0.23\linewidth}L{0.33\linewidth}L{0.36\linewidth}@{}}
\toprule
机制 & 正常路径 & 必须补的反例 \\
\midrule
系统调用 & 参数合法，返回字节数或结果 & 用户指针跨页且后一页不可读，返回 \code{-EFAULT} 或部分成功 \\
调度 & runnable 任务被选中运行 & 任务睡在 wait queue 中，不能因为优先级高就运行 \\
锁 & 持锁者修改共享状态后释放 & 两把锁顺序反过来，形成 wait-for 环 \\
缺页 & VMA 合法，分配或读入页面 & 地址不属于任何 VMA，必须发信号而不是分配页 \\
DMA & completion 到达后回收 buffer & completion 前释放 buffer，设备可能写坏别人的内存 \\
持久化 & 日志提交后恢复可重放 & 写了数据但没写 commit，崩溃后不能当成功 \\
\bottomrule
\end{longtable}

一个机制是否成立，先看没有它时哪里坏，再看新增了什么状态对象，最后看正常路径和反例能不能守住边界。

\begin{longtable}{@{}L{0.18\linewidth}L{0.39\linewidth}L{0.35\linewidth}@{}}
\toprule
层次 & 必须回答的问题 & 例子 \\
\midrule
原始失败 & 没有这个机制时，哪个具体动作会坏 & 死循环霸占 CPU；用户指针越界；写一半掉电 \\
状态对象 & 内核新增哪本账，记录哪些字段 & task state、runqueue、VMA、PTE、request、journal record \\
状态转移 & 谁能改账本，入口是什么 & timer interrupt 改调度状态；page fault 安装 PTE；completion 完成 request \\
不变量 & 哪些关系不能被破坏 & sleeping 任务不能在 runqueue；DMA 完成前 buffer 不能释放 \\
反例测试 & 怎样故意触发边界 & lost wakeup、非法地址、共享页写入、I/O error、崩溃恢复 \\
\bottomrule
\end{longtable}

例如调度不能只列 FCFS、RR、CFS，而要回答：死循环为什么能霸占 CPU；定时器如何给内核抢占入口；task 在 runnable、running、blocked 之间怎样转移；runqueue 用队列、树或多队列时复杂度怎么变；lost wakeup 和饥饿怎样测出来。虚拟内存也一样：直接物理地址会互相破坏，VMA 表达地址区间意图，PTE 记录单页翻译事实，page fault 把硬件证据交给内核决策，非法地址、共享页写入、权限错误和 TLB 失效分别压在不同边界上。

保存路径最终会落到五个维度。

\begin{longtable}{p{0.18\linewidth}p{0.34\linewidth}p{0.36\linewidth}}
\toprule
维度 & 追问 & 例子 \\
\midrule
身份 & 对象如何被引用 & pid、tid、fd、inode number、socket 四元组、DMA tag \\
权限 & 谁能做什么 & 页表读写执行位、fd 标志、资源能力、命名空间 \\
生命周期 & 何时创建和释放 & fork/exit/wait、open/dup/close、page refcount、request completion \\
等待 & 无法立即推进时停在哪里 & runqueue、wait queue、page fault、I/O queue、timer \\
提交 & 何时对外可见或持久 & 系统调用返回、PTE 安装、DMA 完成、fsync、提交记录 \\
\bottomrule
\end{longtable}

只说其中一个维度，很容易把机制讲偏。比如“socket 可读”只覆盖了等待维度，还要继续追：哪个 fd 引用它，谁有权限读，缓冲区里数据生命周期如何，读取后队列怎样变化，关闭时如何唤醒等待者。

这些维度还要放回层次关系中。操作系统经常呈现层次结构：硬件在底层，内核抽象在中间，用户 API 在上层。但真实问题会反向穿透层次。一个 \code{write} 系统调用看似属于文件 API，却可能等待内存回收、块层队列、驱动中断、设备 flush 和调度器唤醒。一个 page fault 看似属于内存管理，却可能读取文件页，触发块设备 I/O。一个安全拒绝看似属于权限系统，却依赖路径解析、命名空间和当前进程的身份信息。

\begin{lstlisting}
user write()
  -> fd table
  -> open file object
  -> page cache
  -> writeback
  -> block layer
  -> driver DMA
  -> interrupt completion
  -> wakeup
\end{lstlisting}

真实系统不是互不相干的模块，而是交叉依赖的状态机网络。

保存请求不能只是一段普通函数代码，因为它至少要穿过五个关口。

\begin{longtable}{p{0.18\linewidth}p{0.36\linewidth}p{0.34\linewidth}}
\toprule
关口 & 要解决的问题 & 保存案例中的落点 \\
\midrule
机器能保存状态 & CPU、寄存器、栈、时钟和设备寄存器如何形成可恢复现场 & task 能从睡眠处回来，设备 ready 位能被观察 \\
机器能被打断 & 中断、异常和定时器怎样把外部事件交给内核 & 设备完成能唤醒等待者，死循环能被抢占 \\
地址能被检查 & MMU、页表、VMA 和 PTE 怎样区分合法缺页和非法访问 & 用户 buffer 跨页时，内核知道复制、补页还是拒绝 \\
资源能被引用 & fd、打开文件对象、inode、socket、request 怎样把一次操作变成对象账本 & 整数 fd 找到文件对象，I/O completion 找回请求 \\
结果能被证明 & page cache、writeback、journal、fsync 和 rename 怎样推进提交边界 & “保存成功”从函数返回推进到崩溃后可恢复 \\
\bottomrule
\end{longtable}

这些关口也可以压成一张“失败、机制、代价”表：

\begin{longtable}{p{0.24\linewidth}p{0.34\linewidth}p{0.3\linewidth}}
\toprule
失败 & 需要的机制 & 新增代价 \\
\midrule
固定电路不通用 & 存储程序 CPU & 指令编码、取指译码、异常语义 \\
单核状态无法被安全共享 & 多核一致性和原子操作 & cache line 迁移、内存序、锁竞争 \\
程序直接用物理地址 & MMU 和虚拟地址 & 页表、TLB、缺页成本 \\
CPU 搬运所有 I/O 数据 & DMA 和描述符队列 & buffer 生命周期、IOMMU、completion \\
轮询浪费 CPU & 中断 & 异步共享状态 \\
任务死循环霸占 CPU & 定时器抢占 & 上下文保存和调度开销 \\
任务互相覆盖内存 & 地址空间和权限 & 页表、缺页、TLB 成本 \\
用户能随便改设备 & 系统调用和特权级 & 入口检查和复制成本 \\
I/O 等待拖住计算 & 阻塞队列、异步和 completion & 请求身份和 buffer 生命周期 \\
写一半崩溃 & 日志、fsync、提交记录 & 写放大和恢复复杂度 \\
\bottomrule
\end{longtable}

OS 测试不是只看程序有没有跑完，而是验证契约有没有守住。正常路径、错误路径和边界路径都要能指出是哪条不变量断了。

\begin{rubricbox}
\begin{itemize}
\tightlist
\item 系统调用能拆成用户态入口、内核对象、权限检查、状态更新和返回值。
\item “函数返回成功”“内核接受请求”“设备完成”“业务提交生效”不能混成同一个完成点。
\item 等待现象必须落到具体队列：CPU、锁、页面、I/O、子进程都不是同一种等待。
\item 隔离边界必须落到硬件或内核对象：地址空间、文件权限、用户/内核态、命名空间各管不同问题。
\item 任意机制都应有状态、入口、失败和证据四列。
\end{itemize}
\end{rubricbox}

反例同样重要。若仍然把 fd 当成文件本体，把 \code{write} 成功当成落盘，把线程 ready 当成 running，把 \code{malloc} 成功当成物理页已分配，把 \code{notify\_one} 当成消息投递，就说明 OS 契约还没有建立。

\section{综合验证：从机制到证据}

本章每个关键断言都可以在一台真实 Linux 机器上复查。下面这组实验不需要特殊环境；本章展示的输出就是在本书实验机（aarch64 容器）上跑出来的，你的数字可能不同，但边界关系应当一致。

\begin{enumerate}
\tightlist
\item 坏地址实验：把 \code{write} 的 buffer 指向 \code{0x1}，观察程序打印 \code{n=-1 errno=14 (Bad address)}；再用 \code{strace -e trace=write} 对照系统调用边界上的真实一行。
\item offset 归属实验：分别跑 \code{dup}、重新 \code{open}、\code{fork} 三个版本，用 \code{od -An -c data} 验证最终内容是 \code{CB}、\code{B}、\code{CP}，并解释差异来自打开文件对象的共享关系。
\item fd 观察：让 \code{fd\_demo} 睡眠，读 \code{/proc/<pid>/fd} 和 \code{/proc/<pid>/status}，确认两个 fd 指向同一个文件、进程状态是 \code{S (sleeping)}。
\item 完成点距离实测：给 64KiB 的 \code{write} 和紧随其后的 \code{fsync} 分别计时，比较两者差几个数量级；换存储介质（tmpfs、SSD、机械盘）再测一次，观察差距怎样变化。
\item 崩溃点推演：在完成点时间线上任选一点假设“断电”，写出重启后用户、内核、设备各自能看到什么，并说明哪一层的账本足以支撑恢复结论。
\end{enumerate}

做到这五条，本章的契约就不是背下来的，而是被自己亲手推翻过又重建的。

操作系统不是一堆彼此孤立的子系统，而是在维护一组可验证契约。遇到任何机制，都要先找没有它时哪个具体动作会坏，再问内核新增了哪本账，谁能通过哪个入口改这本账，不能继续时睡在哪里，成功和失败分别留下什么证据。保存按钮只是一个入口；同样的方法也适用于调度、虚拟内存、文件系统、网络和设备驱动。
